\chapter{Universal homeomorphisms in positive characteristic}
\index{universal homeomorphism!envelope theorem}\index{positive characteristic!universal homeomorphism}
\label{chap:bj-universal-homeomorphisms}

\section{Orientation}
\label{sec:ch43-orientation}

\subsection{What is proved here, and out of what}
\label{sec:ch43-what}

The normal theorem of \Cref{chap:bj-alterations} is not yet the desired
result for a nonnormal unibranch variety.
Its normalization is finite and birational, but the finite descent theorem of
\Cref{chap:bj-finite-descent} requires a normal target, and it descends the
envelope property rather than psh functions themselves.
The correct replacement is stronger and more pointwise:
pullback through the normalization will be shown to identify the two psh
spaces.

This chapter establishes the topological part of that statement in every
characteristic and proves psh reflection in characteristic \(p>0\).
The coefficient descent is especially clean in positive characteristic:
one Frobenius power moves every local coefficient downstairs.
\Cref{chap:bj-subintegral-normalization} constructs the
characteristic-zero replacement.

Let
\[
                         \nu:Y\longrightarrow X
\tag{43.1}\label{eq:bj-universal-map}
\]
be a finite birational universal homeomorphism of integral projective
varieties over the algebraically closed trivially valued field \(k\).
We assume that \(Y\) is normal.
Neither the topology nor the coefficient argument below uses the envelope
property.

The chapter has exactly two mathematical contents, and they are of quite
different natures.  The first is topological and characteristic-free:
\(\nu^{\An}\) is a homeomorphism, so a function on \(X^{\An}\) and its
pullback to \(Y^{\An}\) are literally the same function of the same points,
including at semivaluations whose support is a proper closed subvariety and
including where the value is \(-\infty\)
(\Cref{lem:bj-universal-analytic-homeomorphism}).  Nothing about psh
functions is involved; what is involved is the uniqueness of an absolute
value on a purely inseparable residue extension, and the compactness of the
analytification of a projective variety.

The second content is algebraic and is where the characteristic decides
everything.  Being the same function is not the same as being psh in the two
senses, because the two senses are defined by different families of
approximants: a generically finite Fubini--Study function upstairs is built
from sections of a line bundle on \(Y\), whose local coefficients live in
\(\cO_Y\) and need not lie in \(\cO_X\).  The task is to convert an upstairs
GF presentation into a downstairs one \emph{without changing the function},
and in particular without averaging over fibres --- \(\nu\) has singleton
fibres, so there is nothing to average, and division by a fibre degree would
be division by one.  In characteristic \(p\) the conversion is performed by
raising every local coefficient to a single Frobenius power
(\Cref{lem:bj-uniform-frobenius}, \Cref{prop:bj-frobenius-gf-descent}); in
characteristic zero there is no such operation, and the conversion is
performed instead by the additive-power theorem of Roberts and Singh, in
\Cref{chap:bj-subintegral-normalization}.

\subsection{The uniformity, which is the whole point}
\label{sec:ch43-uniformity}

It is worth being precise, before any proof, about what the
characteristic-\(p\) argument actually asserts, because the summary
``Frobenius'' conceals the only step that requires an idea.

The Stacks Project criterion for an affine universal homeomorphism in
characteristic \(p\) is \emph{elementwise}: given \(b\in B\), some power
\(b^{p^{e(b)}}\) lies in \(A\).  Taken literally this is useless for
descending a Fubini--Study function.  A GF function is
\[
 \psi=\frac1m\max_j\{\log|s_j|+\lambda_j\},
 \qquad
 s_j\in H^0\bigl(Y,\nu^*(mL)\bigr),
\]
and its descent requires the local coefficients \(f_{U,j}\) of the
\emph{finitely many} sections \(s_j\), read in frames on the members
\(U\) of \emph{some} affine cover trivializing \(mL\), to be raised to
\emph{one and the same} exponent \(q\).  Three separate uniformities are
needed, and they are logically independent:

\begin{enumerate}
\item\label{itm:ch43-unif:sections}
      one exponent for the finitely many coefficients on a fixed chart;
\item\label{itm:ch43-unif:charts}
      one exponent for all charts of the cover at once;
\item\label{itm:ch43-unif:covers}
      one exponent valid on \emph{every} affine open of \(X\), so that the
      exponent may be chosen \emph{before} \(L\), \(m\) and the trivializing
      cover are known.
\end{enumerate}

Item~\cref{itm:ch43-unif:sections} is bookkeeping: take the maximum of
finitely many exponents and use that \(b^{q}\in A\) implies
\(b^{qp}=(b^q)^p\in A\).  Item~\cref{itm:ch43-unif:charts} is the same
bookkeeping applied to a finite cover, and uses that
\(\{b\in B:b^q\in A\}\) is an \(A\)-subalgebra of \(B\) --- which is true
because \(x\mapsto x^q\) is a ring homomorphism in characteristic \(p\), and
is false in characteristic zero, where the set is not even closed under
addition.

Item~\cref{itm:ch43-unif:covers} is the one that matters, and it is the
reason \Cref{lem:bj-uniform-frobenius} must be supplemented by
\Cref{cor:ch43-frobenius-all-affines}.  The exponent has to be produced once
and for all from the morphism \(\nu\) alone, because the cover trivializing
\(mL\) depends on \(L\) and on \(m\), and both vary along a regularizing net:
the net of \Cref{def:bj-psh}\,(3) has \emph{varying} classes \(L_i\), each
requiring its own cover and its own denominator \(m_i\).  An exponent
attached to one fixed cover would have to be re-chosen at every stage, and
the descended sections at different stages would live in unrelated Veronese
multiples \(q_im_iL_i\) with \(q_i\to\infty\), so that no single downstairs
Veronese multiple accommodates the whole net.  What rescues the argument is that the
characteristic-\(p\) condition \(b^q\in A\) \emph{localizes}: it passes from
a chart to its principal opens and hence, by a sheaf argument, to every
affine open of \(X\).  A single \(q=p^e\), read off from a single finite
cover, therefore serves every chart of every cover simultaneously, forever.

\begin{warning}{Uniformity is not automatic, and the chapter turns on it}
\label{warn:ch43-uniformity}
Three misreadings should be excluded at the outset.

\emph{The exponent is not universal over morphisms.}  There is no exponent
serving all finite universal homeomorphisms at once: for the extensions
\(k[t^{p^n},t^{p^n+1}]\subset k[t]\) the least admissible exponent is
\(p^n\), which is unbounded in \(n\)
(\Cref{ex:ch43-no-uniform-exponent}).  The exponent is uniform \emph{for a
fixed \(\nu\)}, and its existence uses that \(X\) is quasi-compact and
\(\nu\) finite.

\emph{The exponent is not a bound on the degree of anything.}  It is an
exponent of an inseparability, not a degree; the map \(\nu\) may be
birational, so that all its residue extensions at generic points are
trivial, while the exponent is arbitrarily large.

\emph{One exponent per chart is not enough.}  The \(q\)-th powers of local
coefficients transform under a change of frame by \(g^{-q}\), which is the
transition cocycle of \(qmL\).  If two overlapping charts used exponents
\(q_1\ne q_2\), the resulting local functions would be coefficients of
sections of two \emph{different} line bundles and would not glue.  The fix
is to raise everything to a common exponent, which is legitimate but
requires the common exponent to exist
(\Cref{ex:ch43-two-cusps}).
\end{warning}

The contrast with \Cref{chap:bj-subintegral-normalization} is the cleanest
illustration in \cref{part:bj} of how much characteristic \(p\) can
simplify.  There, in characteristic zero, the extension
\(\cO_X\subset\nu_*\cO_Y\) is subintegral, and the substitute for ``raise to
the power \(q\)'' is the Roberts--Singh identity: there are an integer
\(c\ge0\), a threshold \(N\), and \emph{fixed} elements
\(c_1,\ldots,c_c\in B\) such that
\[
 b^n+\sum_{i=1}^{c}\binom nic_ib^{n-i}\in A
 \qquad(n\ge N).
\]
The left-hand side is not a power of \(b\), it is a power \emph{corrected by
lower-order terms}, and consequently:
the descended object is not a section but a \emph{system} of sections;
the identity \(\log|b^q|=q\log|b|\) is replaced by a two-sided valuative
sandwich with an additive error \(C\); the resulting functions are not equal
to the original but only bracket it, so a limit \(n\to\infty\) must be
taken; and monotonicity of the resulting net has to be arranged by hand.
Each of those four replacements costs a section of
\Cref{chap:bj-subintegral-normalization}.  In characteristic \(p\) all four
collapse, because \(x\mapsto x^q\) is a ring homomorphism and
\(\log|b^q|=q\log|b|\) holds exactly, at every semivaluation, including where
both sides are \(-\infty\).  That is the entire content of
\Cref{prop:bj-frobenius-gf-descent}, and it is why the present chapter is
the shorter of the two.

\subsection{The firewall}
\label{sec:ch43-firewall}

\begin{warning}{The envelope theorem is not available in this chapter}
\label{warn:ch43-firewall}
Nothing in \Cref{chap:bj-universal-homeomorphisms} may use, directly or
transitively:
\begin{enumerate}
\item \Cref{thm:bj-final-envelope}, in any of its formulations, nor
      \Cref{def:bj-envelope-property-final}, nor
      \Cref{cor:bj-historical-conjecture};
\item the envelope property \(\operatorname{EP}\) of
      \Cref{def:bj-envelope-property} as a \emph{theorem} for any class on
      any variety --- in this chapter \(\operatorname{EP}\) never occurs at
      all except in the single conditional statement
      \Cref{cor:ch43-ep-transport}, where it is displayed as a hypothesis on
      the source \(Y\);
\item closure of any class \(\PSH_X(\theta)\) under suprema, regularized or
      not;
\item any statement of \Cref{chap:bj-subintegral-normalization} or of
      \Cref{chap:bj-envelope-consequences}.  Forward cross-references to
      those chapters occur in remarks, to say where a companion result or an
      assembly is carried out; none of them is used in a proof, and every
      proof below can be read with those remarks deleted.
\end{enumerate}
\Cref{thm:bj-positive-universal-reflection} is an \emph{input} to
\Cref{thm:bj-final-envelope}, not a consequence of it, and the same is true
of every numbered statement below.
\end{warning}

\begin{remark}[Why the firewall is not vacuous here]
\label{rem:ch43-firewall-not-vacuous}
Two places are tempting, and both would be circular.

First, \Cref{thm:bj-positive-universal-reflection} asserts that
\(\nu^*\) is a \emph{bijection}.  Surjectivity of pullback along a
birational morphism of normal projective varieties is available in this book
as \Cref{cor:ch42-birational-bijection}, and along a finite morphism as
\Cref{cor:ch41-finite-bijection}.  Neither applies: \(X\) is not normal ---
if it were, \(\nu\) would be an isomorphism and there would be nothing to
prove --- and both of those corollaries are themselves consequences of the
normal envelope theorem, hence unavailable for producing an input to the
unibranch one.  The bijection proved here is obtained from scratch, by
converting presentations.

Second, at the last step of the proof of
\Cref{thm:bj-positive-universal-reflection} one has a decreasing net of
\(L_i\)-psh functions on \(X\) with \(c_1(L_i)\to\theta\) and wants their
limit to be \(\theta\)-psh.  This is \Cref{thm:bj-basic-closure}\,(5), the
closure of \(\PSH\) under \emph{decreasing} nets with varying classes, which
is part of the definitional package of \Cref{chap:bj-fs-pl-psh} and is
independent of every envelope statement.  It is not the closure under
\emph{increasing} families, which is the envelope property and is exactly
what may not be used.  The two are distinguished at every appeal below, and
the one exceptional alternative in the decreasing statement --- an
identically \(-\infty\) limit --- is excluded each time by recording the
value at \(v_\triv\).
\end{remark}

\subsection{The five tasks, and where each is discharged}
\label{sec:ch43-tasks}

\begin{enumerate}
\item\label{itm:ch43-task:radicial}
      \emph{Radicial finite maps and topological invariance.}  Assemble the
      three Stacks Project criteria into the exact statements used
      (\Cref{sec:ch43-radicial}), and prove that \(\nu^{\An}\) is a
      homeomorphism carrying supports to supports and centers to centers,
      with the consequence that upper semicontinuous regularization and
      pointwise suprema commute with pullback
      (\Cref{sec:ch43-analytic}).
\item\label{itm:ch43-task:class}
      \emph{Promotion of the numerical class.}  A regularizing net upstairs
      has classes that merely converge to \(\nu^*\theta\); they must be
      replaced, without changing a single value of a single function, by
      classes that are honest pullbacks \(\nu^*L_i\)
      (\Cref{sec:ch43-promotion}).  This is the step that makes the frame
      calculation of task~\cref{itm:ch43-task:reflect} legitimate.
\item\label{itm:ch43-task:frobenius}
      \emph{Uniform Frobenius powers.}  Produce a single \(q=p^e\), depending
      on \(\nu\) alone, with \(b^q\in\cO_X(U)\) for every affine open
      \(U\subseteq X\) and every \(b\in\cO_Y(\nu^{-1}U)\)
      (\Cref{sec:ch43-frobenius}).  The passage from a fixed finite cover to
      all affine opens is \Cref{cor:ch43-frobenius-all-affines} and is not
      cosmetic.
\item\label{itm:ch43-task:reflect}
      \emph{Pullback and reflection of Fubini--Study functions.}  Show that
      \(\nu^*\) is a bijection from \(\cH^{\mathrm{gf}}_{\mathbf Q}(X,L)\)
      onto \(\cH^{\mathrm{gf}}_{\mathbf Q}(Y,\nu^*L)\), preserving the
      \(-\infty\)-locus (\Cref{sec:ch43-gf-descent}).
\item\label{itm:ch43-task:psh}
      \emph{Psh reflection and descent through normalization.}  Combine the
      previous three into
      \Cref{thm:bj-positive-universal-reflection}, and identify the
      normalization of a unibranch variety as a map to which it applies
      (\Cref{sec:ch43-reflection}, \Cref{sec:ch43-unibranch}).
\end{enumerate}

\subsection{Standing hypotheses and notation}
\label{sec:ch43-hypotheses}

Throughout, \(k\) is an algebraically closed field carrying the trivial
absolute value, and \(\nu:Y\to X\) is as in \eqref{eq:bj-universal-map}.
\Cref{sec:ch43-radicial,sec:ch43-promotion} and
\Cref{sec:ch43-unibranch} are valid in every characteristic;
\Crefrange{sec:ch43-frobenius}{sec:ch43-reflection} assume
\(\operatorname{char}k=p>0\) and say so at their head.

Numerical classes are real unless stated otherwise, and
\(N^1(X)=N^1(X)_{\mathbf R}\).  A rational line bundle is an element of
\(\Pic(X)_{\mathbf Q}\).  The analytic vocabulary is that of
\Cref{chap:trivial-analytification}: for \(v\in X^{\An}\) we write \(s(v)\)
for its support (\Cref{def:triv-support}) and \(c(v)\) for its center, and a
point with \(s(v)\subsetneq X\) is said to have proper support
(\Cref{def:ch32-proper-support}).  Sections, ideals and their logarithms are
normalized as in \Cref{chap:bj-fs-pl-psh}: \(\log|s|(v)=-v(s)\), computed in
any local frame at \(c(v)\), with the value \(-\infty\) permitted and never
discarded.  The sign conventions of \Cref{app:normalizations} are in force
without exception.

We write \(\cH^{\mathrm{gf}}_{\mathbf Q}(L)\) for the rational generically
finite Fubini--Study functions attached to \(L\in\Pic(X)_{\mathbf Q}\)
(\Cref{def:gf-fs-functions}), and \(\PSH_X(\theta)\) for the \(\theta\)-psh
functions in the sense of \Cref{def:bj-psh}; recall that the approximating
net in that definition has \emph{varying} classes, a point which
\Cref{sec:ch43-promotion} exists to exploit.

Finally, for an affine open \(U\subseteq X\) we abbreviate
\[
 A_U:=\cO_X(U),
 \qquad
 B_U:=\cO_Y(\nu^{-1}U),
\]
noting that \(\nu^{-1}U\) is affine because \(\nu\) is finite, hence affine,
and that \(A_U\to B_U\) is an injective finite ring map because \(\nu\) is
finite, birational, and dominant between integral schemes.

\subsection{What is used from earlier chapters}
\label{sec:ch43-dependencies}

The imports are few and are listed here so that the firewall can be checked
by inspection.

From \Cref{chap:trivial-analytification}: the description of \(X^{\An}\) as
pairs (scheme point, absolute value on its residue field)
(\Cref{lem:ch41-point-pair}); compactness and Hausdorffness of \(X^{\An}\)
for \(X\) projective (\Cref{thm:triv-berkovich-compactness}); the
support and center maps and their (anti)continuity
(\Cref{thm:ch32-support-center-maps}); and the scaling action
(\Cref{prop:ch32-scaling-action}).

From \Cref{chap:bj-fs-pl-psh}: the definition of
\(\cH^{\mathrm{gf}}_\Lambda(L)\) and its elementary calculus
(\Cref{def:gf-fs-functions}, \Cref{prop:ch34-fs-operations}); the definition
of \(\PSH_X(\theta)\) (\Cref{def:bj-psh}) with varying classes; and the
closure package \Cref{thm:bj-basic-closure}, of which only items~(1),~(2)
and~(5) are used.

From \Cref{chap:bj-finite-descent}: the analytic calculus of finite
morphisms, namely \Cref{lem:ch41-finite-surjective},
\Cref{lem:ch41-support-centre-finite}, and the radicial homeomorphism
criterion \Cref{lem:ch41-radicial-homeo}, of which
\Cref{lem:bj-universal-analytic-homeomorphism} is the version adapted to a
possibly nonnormal target.  Nothing from the averaging machinery of that
chapter is used: there is no determinant norm, no flattening, and no
division by a degree anywhere below.

From \Cref{chap:bj-alterations}: nothing is used in the proofs.  That
chapter supplies the theorem that will be combined with the present one in
\Cref{chap:bj-envelope-consequences}, and the combination is performed
there, not here.

From \Cref{app:valuation-algebra} and \Cref{app:stacks}: the definitions of
radicial morphism and universal homeomorphism
(\Cref{def:appA-universal-homeomorphism}), the criterion
\Cref{thm:appA-universal-homeomorphism-criterion}, and the Stacks Project
tags 01S4, 04DF, 0CNF and 035S, each audited in
\Cref{sec:stacks-universal-homeomorphisms} with its hypotheses and its first
place of use.  No tag is quoted below that is absent from that audit.

\subsection{Roadmap}
\label{sec:ch43-roadmap}

\Cref{sec:ch43-radicial} collects the criteria for radiciality and universal
homeomorphy and illustrates them.  \Cref{sec:ch43-analytic} proves the
analytic homeomorphism and its consequences for regularization and suprema.
\Cref{sec:ch43-promotion} promotes the numerical class of a regularizing
net.  \Cref{sec:ch43-frobenius} constructs the uniform Frobenius exponent
and localizes it to all affine opens; this is the heart of the chapter.
\Cref{sec:ch43-gf-descent} descends a single GF stage, and
\Cref{sec:ch43-reflection} assembles the reflection theorem.
\Cref{sec:ch43-unibranch} identifies the normalization of a unibranch
variety as an instance.  \Cref{sec:ch43-computations} carries four examples
to explicit numbers, \Cref{sec:ch43-audit} performs the extended-real and
circularity audit, \Cref{sec:ch43-failure} maps the failure boundaries, and
\Cref{sec:ch43-dictionary} and \Cref{sec:ch43-source-note} close the chapter.

\section{Radicial maps and universal homeomorphisms}
\label{sec:ch43-radicial}

\begin{definition}[Radicial and universal homeomorphism]
\label{def:bj-universal-homeomorphism}
A morphism is \emph{radicial} if it is universally injective.
A morphism is a \emph{universal homeomorphism} if every base change induces a
homeomorphism on underlying topological spaces.
\end{definition}

\subsection{The three criteria, stated exactly}
\label{sec:ch43-criteria}

The criteria used here are exact.
A morphism is radicial if and only if its diagonal is surjective, by
\cite[\href{https://stacks.math.columbia.edu/tag/01S4}{Tag 01S4}]
{stacks-project}.
A morphism is a universal homeomorphism if and only if it is integral,
surjective, and radicial, by
\cite[\href{https://stacks.math.columbia.edu/tag/04DF}{Tag 04DF}]
{stacks-project}.
For the finite map \(\nu\), integrality is automatic.

We record the three criteria in the form in which they will be applied,
together with the fourth --- the characteristic-\(p\) affine criterion ---
which is the engine of \Cref{sec:ch43-frobenius}.

\begin{proposition}
\label{prop:ch43-criteria}
Let \(f:W\to Z\) be a morphism of schemes.
\begin{enumerate}
\item\label{itm:ch43-crit:radicial}
      The following are equivalent: \(f\) is universally injective;
      \(f\) is radicial, that is, injective on points and inducing purely
      inseparable residue-field extensions; the diagonal
      \(\Delta_f:W\to W\times_ZW\) is surjective; for every field \(K\), the
      map \(W(K)\to Z(K)\) is injective.
\item\label{itm:ch43-crit:uh}
      \(f\) is a universal homeomorphism if and only if it is integral,
      surjective and universally injective.
\item\label{itm:ch43-crit:finite}
      If \(f\) is of finite type and a universal homeomorphism, then \(f\)
      is finite.
\item\label{itm:ch43-crit:charp}
      Suppose \(Z=\Spec A\) and \(W=\Spec B\) are affine and that \(A\) and
      \(B\) are \(\mathbf F_p\)-algebras.  Then \(f\) is a universal
      homeomorphism if and only if the kernel of \(A\to B\) is locally
      nilpotent and every \(b\in B\) admits a \(p\)-power \(q\) with
      \(b^q\) in the image of \(A\).
\end{enumerate}
\end{proposition}

\begin{proof}[Foundational input]
Items~\cref{itm:ch43-crit:radicial} and~\cref{itm:ch43-crit:uh} are
\cite[\href{https://stacks.math.columbia.edu/tag/01S4}{Tag 01S4}]
{stacks-project} and
\cite[\href{https://stacks.math.columbia.edu/tag/04DF}{Tag 04DF}]
{stacks-project}, imported as
\Cref{def:appA-universal-homeomorphism} and
\Cref{thm:appA-universal-homeomorphism-criterion} and audited in
\Cref{sec:stacks-universal-homeomorphisms}.
Item~\cref{itm:ch43-crit:finite} is the combination of
\cref{itm:ch43-crit:uh} with the standard fact that a morphism which is
both integral and of finite type is finite; we use it only to observe that
in the setting of \eqref{eq:bj-universal-map} the hypothesis ``finite'' is
not an extra assumption once ``universal homeomorphism'' and ``of finite
type'' are granted.
Item~\cref{itm:ch43-crit:charp} is
\cite[\href{https://stacks.math.columbia.edu/tag/0CNF}{Tag 0CNF}]
{stacks-project}; that tag is stated for a ring map inducing an isomorphism
\(A[1/p]\to B[1/p]\) and requires that both \(qb\) and \(b^q\) lie in the
image.  When \(A\) and \(B\) are \(\mathbf F_p\)-algebras one has
\(A[1/p]=0=B[1/p]\) and \(qb=0\), so the hypothesis is vacuous and the first
condition drops out.  This is the reading recorded in
\Cref{sec:stacks-universal-homeomorphisms}.
\end{proof}

\begin{remark}[Which direction of each criterion is used]
\label{rem:ch43-criteria-direction}
The bookkeeping matters, because the two directions have different costs.

For \(\nu\) of \eqref{eq:bj-universal-map} we are \emph{given} that it is a
universal homeomorphism and \emph{use} radiciality: this is the easy
direction of \cref{itm:ch43-crit:uh}, and it is what feeds
\Cref{lem:bj-universal-analytic-homeomorphism}.  For the normalization of a
unibranch variety, in \Cref{lem:bj-unibranch-normalization-universal}, the
implication is used in the other direction: bijectivity is known, radiciality
is deduced from the diagonal criterion \cref{itm:ch43-crit:radicial}, and
\cref{itm:ch43-crit:uh} then produces the universal homeomorphism.  In
\Cref{sec:ch43-frobenius} the criterion \cref{itm:ch43-crit:charp} is used
in the direction ``universal homeomorphism \(\Rightarrow\) elementwise
\(p\)-powers''; the converse direction is never needed.
\end{remark}

\begin{warning}{Bijective on closed points is strictly weaker}
\label{warn:ch43-bijective-not-enough}
Over a non-algebraically-closed field a finite morphism can be bijective on
closed points without being radicial: the residue extensions may be
separable.  For \(k=\mathbf F_p(s)\), the map
\(\Spec k(s^{1/p})\to\Spec k\) is finite, bijective on points, and a
homeomorphism; here it \emph{is} radicial, since the residue extension is
purely inseparable.  If \(p\ne2\), the equally one-point map
\(\Spec k(\sqrt s)\to\Spec k\) is finite and bijective on points, but its
residue extension is separable of degree two, the diagonal is not
surjective, and after base change to \(k(\sqrt s)\) the source splits.  It
is therefore not a universal homeomorphism.  This is precisely why
\Cref{lem:bj-unibranch-normalization-universal} is stated over an
algebraically closed field and why its proof passes through the diagonal
rather than through bijectivity alone; see
\Cref{rem:ch43-where-algebraically-closed}.
\end{warning}

\subsection{Stability properties}
\label{sec:ch43-stability}

\begin{lemma}
\label{lem:ch43-uh-stability}
\begin{enumerate}
\item\label{itm:ch43-stab:comp}
      A composition of universal homeomorphisms is a universal
      homeomorphism, and if \(g\circ f\) and \(g\) are universal
      homeomorphisms with \(g\) separated, so is \(f\).
\item\label{itm:ch43-stab:bc}
      Universal homeomorphisms are stable under arbitrary base change and
      under restriction to opens: if \(f:W\to Z\) is a universal
      homeomorphism and \(U\subseteq Z\) is open, then
      \(f^{-1}U\to U\) is a universal homeomorphism.
\item\label{itm:ch43-stab:red}
      A closed immersion defined by a nilpotent ideal is a universal
      homeomorphism; in particular \(Z_{\mathrm{red}}\to Z\) is one.
\item\label{itm:ch43-stab:frob}
      If \(Z\) is an \(\mathbf F_p\)-scheme, the absolute Frobenius
      \(F_Z:Z\to Z\) is a universal homeomorphism.
\end{enumerate}
\end{lemma}

\begin{proof}
\cref{itm:ch43-stab:comp} Integrality, surjectivity and universal
injectivity are each stable under composition; for the cancellation
statement, \(f\) is separated because \(g\) is and \(g\circ f\) is, it is
integral by cancellation for integral morphisms, it is surjective because
\(g\) is injective and \(g\circ f\) is surjective, and it is universally
injective because \(g\circ f\) is.  Apply
\Cref{prop:ch43-criteria}\,\cref{itm:ch43-crit:uh}.

\cref{itm:ch43-stab:bc} All three properties in
\cref{itm:ch43-crit:uh} are stable under base change, and an open
immersion is a base change.

\cref{itm:ch43-stab:red} Such an immersion is finite, hence integral;
it is surjective and injective on points with equality of residue fields.

\cref{itm:ch43-stab:frob} \(F_Z\) is the identity on the underlying
space, hence a homeomorphism; it is affine and integral because
\(a\mapsto a^p\) makes \(\cO_Z\) into an integral algebra over its image
(each section satisfies \(T^p-a^p=(T-a)^p\)); and it is universally
injective because for any field \(K\) the map \(x\mapsto x^p\) is injective.
Alternatively, apply \cref{itm:ch43-crit:charp} with \(q=p\).
\end{proof}

\begin{example}[Frobenius: a universal homeomorphism that is not birational]
\label{ex:ch43-frobenius-uh}
Let \(k\) be algebraically closed of characteristic \(p>0\) and
\(X=\mathbf P^1_k\).  The \(k\)-linear relative Frobenius
\[
 F:\mathbf P^1\longrightarrow\mathbf P^1,
 \qquad
 F^\#(t)=t^p
\]
is finite of degree \(p\), surjective, and radicial: its fibres are
singletons and the residue extension at the generic point is
\(k(t)\subset k(t^{1/p})\), purely inseparable of degree \(p\).  By
\Cref{prop:ch43-criteria}\,\cref{itm:ch43-crit:uh} it is a universal
homeomorphism.  It is \emph{not} birational, and it is not an isomorphism.

The example separates two hypotheses that the running map
\eqref{eq:bj-universal-map} bundles together.  Universal homeomorphy is a
statement about points and topology; birationality is a statement about
function fields.  \Cref{prop:ch43-frobenius-no-birational} below records
that the characteristic-\(p\) descent argument of this chapter uses only the
former.
\end{example}

\begin{example}[The cuspidal cubic: radicial, birational, not an isomorphism]
\label{ex:ch43-cusp-basic}
Let \(C\subset\mathbf P^2\) be the cuspidal cubic
\(\{ZY^2=X^3\}\), with the singular point \(P=[0:0:1]\).  Its normalization
is
\[
 \nu:\mathbf P^1\longrightarrow C,
 \qquad
 [u:w]\longmapsto[u^2w:u^3:w^3],
\]
which satisfies \(ZY^2=w^3u^6=(u^2w)^3=X^3\); it is a bijection --- the
point \(P\) has the unique preimage \([0:1]\) --- and induces an isomorphism
on all residue fields, all of which are \(k\).  On the affine chart
\(w\ne0\), writing \(t=u/w\) and \(x=X/Z=t^2\), \(y=Y/Z=t^3\), it is the
inclusion
\[
 A=k[t^2,t^3]\subset B=k[t].
\]
It is finite, birational, injective on points, surjective, and a
homeomorphism for the Zariski topology; by
\Cref{prop:ch43-criteria}\,\cref{itm:ch43-crit:uh} it is a universal
homeomorphism.  It is not an isomorphism, since \(t\notin A\).

The example is the running one for the whole chapter.  Note already the two
faces it presents.  Topologically \(\nu\) discards nothing: \(\Spec B\) and
\(\Spec A\) have the same points, and after analytification they will have
the same Berkovich points as well.  Algebraically \(\nu\) discards a great
deal: \(B/A\) is a one-dimensional \(k\)-vector space, spanned by the class
of \(t\), and it is exactly that missing \(t\) which prevents a Fubini--Study
function written with the coefficient \(t\) from being visibly a
Fubini--Study function downstairs.  Raising to the power \(p\) repairs the
defect, because \(t^p\in A\) for every prime \(p\); raising to no power
repairs it in characteristic zero, which is the entire difficulty of
\Cref{chap:bj-subintegral-normalization}.
\end{example}

\begin{example}[A universal homeomorphism with nonreduced source]
\label{ex:ch43-nilpotent-uh}
Let \(Z\) be any scheme and \(W\subset Z\) the closed subscheme defined by a
nilpotent ideal sheaf.  Then \(W\to Z\) is a universal homeomorphism by
\Cref{lem:ch43-uh-stability}\,\cref{itm:ch43-stab:red}, in every
characteristic.  The example shows that ``universal homeomorphism'' by
itself carries no information about characteristic, and that the standing
hypothesis in \eqref{eq:bj-universal-map} that \(X\) and \(Y\) be
\emph{integral} is doing work: it is what forces \(A_U\to B_U\) to be
injective and makes the descent problem one about a genuine inclusion of
rings rather than about nilpotents.
\end{example}

\subsection{Analytic points}
\label{sec:ch43-analytic-points}

An analytic point of \(X^{\An}\) consists of a support prime
\(\mathfrak p\) and a real valuation, or equivalently a multiplicative
seminorm, on the residue field \(\kappa(\mathfrak p)\).
Full-support valuations correspond to \(\mathfrak p=(0)\).
Proper-support semivaluations correspond to nonzero \(\mathfrak p\).
The latter cannot be discarded in a normalization argument.

This is \Cref{lem:ch41-point-pair}, and it is the only description of
\(X^{\An}\) used in this chapter.  The reason to insist on it, rather than
on the more familiar picture of \(X^{\An}\) as a space of valuations on
\(K(X)\), is that the whole chapter is about the interaction of the
morphism \(\nu\) with the \emph{closed} points and with the loci where a GF
function is \(-\infty\).  A description of \(X^{\An}\) that quietly
restricts to full support would make \Cref{lem:ch41-radicial-homeo} and
\Cref{lem:bj-universal-analytic-homeomorphism} appear to be statements about
\(K(X)\) alone, and would silently discard exactly the points at which the
descent is delicate.

\section{The analytic homeomorphism}
\label{sec:ch43-analytic}

\begin{lemma}
\label{lem:ch43-inseparable-extension}
Let \(F\subset E\) be a purely inseparable field extension of characteristic
\(p>0\) and let \(|\cdot|\) be a multiplicative non-Archimedean absolute
value on \(F\).  Then \(|\cdot|\) extends to \(E\) in exactly one way,
namely by
\[
 |b|:=|b^{p^r}|^{1/p^r}
 \qquad\text{whenever }b^{p^r}\in F.
\]
If \(\operatorname{char}k=0\), or more generally if \(F\subset E\) is purely
inseparable and separable, then \(E=F\) and there is nothing to extend.
\end{lemma}

\begin{proof}
\emph{Uniqueness.}  Let \(\lVert\cdot\rVert\) be any multiplicative
extension and let \(b\in E\).  Purely inseparable means that some
\(a:=b^{p^r}\) lies in \(F\).  Multiplicativity gives
\(\lVert b\rVert^{p^r}=\lVert b^{p^r}\rVert=|a|\), and \(\lVert b\rVert\ge0\)
is a real number, so \(\lVert b\rVert=|a|^{1/p^r}\) is determined.  This is
the assertion, in multiplicative notation, that in characteristic \(p\) the
equation \(b^{p^r}=a\) \emph{forces} \(|b|=|a|^{1/p^r}\); nothing is chosen.

\emph{Well-definedness.}  If also \(b^{p^{r'}}=a'\in F\) with \(r'\ge r\),
then \(a'=a^{p^{r'-r}}\), whence
\(|a'|^{1/p^{r'}}=|a|^{p^{r'-r}/p^{r'}}=|a|^{1/p^r}\).  The formula is
therefore unambiguous.

\emph{It is an absolute value.}  Set \(|b|:=|b^{p^r}|^{1/p^r}\).  Then
\(|b|=0\) if and only if \(b=0\), and \(|1|=1\).  For \(b,c\in E\) choose one
\(r\) working for both; multiplicativity of \(|\cdot|\) on \(F\) gives
\[
 |bc|=\bigl|(bc)^{p^r}\bigr|^{1/p^r}
     =\bigl|b^{p^r}c^{p^r}\bigr|^{1/p^r}
     =|b|\,|c|.
\]
For the ultrametric inequality, use the Frobenius identity
\((b+c)^{p^r}=b^{p^r}+c^{p^r}\), valid in characteristic \(p\):
\[
 |b+c|=\bigl|b^{p^r}+c^{p^r}\bigr|^{1/p^r}
      \le\max\bigl\{|b^{p^r}|,|c^{p^r}|\bigr\}^{1/p^r}
      =\max\{|b|,|c|\},
\]
since \(x\mapsto x^{1/p^r}\) is increasing on \(\mathbf R_{\ge0}\).
Finally the restriction of the new absolute value to \(F\) is the old one,
by the case \(r=0\).  In additive notation this is exactly
\[
 \widetilde v(b)=p^{-r}v(a),
 \qquad b^{p^r}=a,
\]
which is \eqref{eq:bj-purely-inseparable-valuation} below.
\end{proof}

\begin{remark}[Why the formula is forced and not merely available]
\label{rem:ch43-forced-formula}
It is worth isolating what makes the characteristic-\(p\) case rigid.  For a
\emph{separable} extension of degree \(d\), an absolute value on \(F\) can
have up to \(d\) distinct extensions to \(E\), and choosing one is a genuine
choice; this is why the finite descent of \Cref{chap:bj-finite-descent}
must average over the extensions rather than pick one.  For a purely
inseparable extension the equation \(b^{p^r}=a\) has \(b\) as its
\emph{unique} root, and the value \(|b|\) is then pinned by
multiplicativity alone.  The rigidity is the same rigidity that will make
\Cref{prop:bj-frobenius-gf-descent} an identity rather than an inequality:
in both places the only mechanism is that \(x\mapsto x^{p^r}\) is an
injective ring homomorphism whose effect on values is multiplication
by \(p^r\).
\end{remark}

\begin{lemma}\label{lem:bj-universal-analytic-homeomorphism}
The morphism \eqref{eq:bj-universal-map} induces a homeomorphism
\[
                         \nu^{\An}:Y^{\An}\xrightarrow{\sim}X^{\An}.
\tag{43.2}\label{eq:bj-universal-analytic-homeomorphism}
\]
\end{lemma}

\begin{proof}
Because \(\nu\) is radicial and surjective, there is a unique prime
\(\mathfrak q\) of \(Y\) above each prime \(\mathfrak p\) of \(X\), and the
extension
\[
                     \kappa(\mathfrak p)\subset\kappa(\mathfrak q)
\]
is purely inseparable.
Let \(v\) be a real valuation of \(\kappa(\mathfrak p)\).
It has a unique extension to \(\kappa(\mathfrak q)\).
Indeed, for \(b\in\kappa(\mathfrak q)\) choose \(r\) with
\(b^{p^r}=a\in\kappa(\mathfrak p)\) in positive characteristic and set
\[
                         \widetilde v(b)=p^{-r}v(a).
\tag{43.3}\label{eq:bj-purely-inseparable-valuation}
\]
This is forced by additivity and is independent of the chosen exponent.
In characteristic zero a purely inseparable field extension is an
isomorphism, so there is nothing to extend.

Thus every analytic point downstairs has exactly one analytic lift and every
analytic point upstairs restricts to it.
The continuous map \(\nu^{\An}\) is bijective.
Analytifications of projective varieties are compact Hausdorff.
A continuous bijection from a compact space to a Hausdorff space is a
homeomorphism, proving \eqref{eq:bj-universal-analytic-homeomorphism}.
\end{proof}

The proof deserves to be run again more slowly, because every one of its
four steps is a place where a hypothesis is spent, and because
\Cref{chap:bj-subintegral-normalization} will reuse the statement verbatim
in characteristic zero.

\emph{Step one: the fibres are singletons with purely inseparable residue
extensions.}  By hypothesis \(\nu\) is a universal homeomorphism, so by
\Cref{prop:ch43-criteria}\,\cref{itm:ch43-crit:uh} it is integral, surjective
and universally injective, and by
\Cref{prop:ch43-criteria}\,\cref{itm:ch43-crit:radicial} universal
injectivity is radiciality.  Surjectivity gives at least one prime of \(Y\)
over each prime of \(X\); injectivity on points gives at most one; and
radiciality gives that \(\kappa(\mathfrak p)\subset\kappa(\mathfrak q)\) is
purely inseparable.  Note that this is a statement about \emph{all} scheme
points, not only closed points and not only generic points.  In the case of
interest, \(\nu\) is birational, so the extension is trivial at the generic
point; it is at the non-generic points, where \(X\) is singular, that the
extension can be nontrivial in characteristic \(p\), and where it is trivial
in characteristic zero.

\emph{Step two: absolute values extend uniquely.}  This is
\Cref{lem:ch43-inseparable-extension}, applied to
\(F=\kappa(\mathfrak p)\subset E=\kappa(\mathfrak q)\).  Existence and
uniqueness are both needed: existence says \(\nu^{\An}\) is surjective,
uniqueness says it is injective.  Neither is automatic for a general finite
morphism; a separable quadratic extension already carries two extensions of
a given absolute value in general, which is why
\Cref{chap:bj-finite-descent} had to average.

\emph{Step three: bijectivity of \(\nu^{\An}\).}  By
\Cref{lem:ch41-point-pair} a point of \(Y^{\An}\) is a pair
\((\eta,\lVert\cdot\rVert)\) with \(\eta\in Y\) a scheme point and
\(\lVert\cdot\rVert\) an absolute value on \(\kappa(\eta)\) trivial on
\(k\), and similarly for \(X^{\An}\); by
\Cref{lem:ch41-support-centre-finite}\,\cref{itm:ch41-sc:supp} the map
\(\nu^{\An}\) is, in this description,
\[
 (\eta,\lVert\cdot\rVert)\longmapsto
 \bigl(\nu(\eta),\lVert\cdot\rVert|_{\kappa(\nu(\eta))}\bigr).
\]
Steps one and two say precisely that this assignment is a bijection: the
first coordinate is bijective on scheme points, and over a fixed point the
second is bijective by unique extension.  Triviality on \(k\) is preserved
in both directions because \(k\) sits inside both residue fields.

\emph{Step four: from bijective to homeomorphic.}  The map \(\nu^{\An}\) is
continuous by functoriality of analytification, \(Y^{\An}\) is compact and
\(X^{\An}\) is Hausdorff by \Cref{thm:triv-berkovich-compactness}, and a
continuous bijection from a compact space to a Hausdorff space is a closed
map, hence a homeomorphism.  Projectivity of \(Y\) is used here and nowhere
else in the lemma; for \(X\) and \(Y\) merely affine of finite type the same
conclusion holds by a direct check on the defining topology, but we shall
not need it.

We henceforth identify \(Y^{\An}\) with \(X^{\An}\) through
\eqref{eq:bj-universal-analytic-homeomorphism}.
Under this identification, a pullback function is literally the same
extended-real function.
In particular, a value \(-\infty\) at a point supported on
\(Z\subsetneq X\) remains \(-\infty\) at the unique point supported on
\(\nu^{-1}(Z)\).

The identification is compatible with all the structure of
\Cref{chap:trivial-analytification} that will be used, and it is worth
recording exactly which structure that is.

\begin{lemma}
\label{lem:ch43-identification-dictionary}
Under \eqref{eq:bj-universal-analytic-homeomorphism}, for
\(w\in Y^{\An}\) with image \(v=\nu^{\An}(w)\):
\begin{enumerate}
\item\label{itm:ch43-id:supp}
      \(s(v)=\nu(s(w))\) and \(s(w)=\nu^{-1}(s(v))\); in particular \(w\) has
      proper support if and only if \(v\) does, and \(w\) lies in
      \(Y^{\val}\) if and only if \(v\) lies in \(X^{\val}\);
\item\label{itm:ch43-id:centre}
      \(c(v)=\nu(c(w))\) and \(c(w)=\nu^{-1}(c(v))\);
\item\label{itm:ch43-id:open}
      for every open \(U\subseteq X\), the identification restricts to a
      homeomorphism \((\nu^{-1}U)^{\An}\to U^{\An}\);
\item\label{itm:ch43-id:scaling}
      the identification is equivariant for the scaling actions of
      \(\mathbf R_{>0}\) of \Cref{prop:ch32-scaling-action}, and carries
      \(w_\triv\) to \(v_\triv\);
\item\label{itm:ch43-id:trivial}
      it carries the trivial semivaluation \(v_{W,\triv}\) of an integral
      closed subvariety \(W\subseteq Y\) to \(v_{\nu(W),\triv}\), and this
      is a bijection between the sets of such points.
\end{enumerate}
\end{lemma}

\begin{proof}
\cref{itm:ch43-id:supp} By \Cref{lem:ch41-point-pair} the support of a point
is the closure of its associated scheme point, and by
\Cref{lem:ch41-support-centre-finite}\,\cref{itm:ch41-sc:supp} the scheme
point of \(v\) is \(\nu(\xi_w)\).  Since \(\nu\) is a homeomorphism of
schemes, it carries closures to closures and
\(\overline{\{\nu(\xi_w)\}}=\nu(\overline{\{\xi_w\}})\), while
\(\nu^{-1}\) of the right-hand side is \(\overline{\{\xi_w\}}\) by
injectivity.  For the last clause, \(s(w)\) is an irreducible component of
\(Y\) --- that is, \(Y\) itself --- if and only if \(\xi_w\) is the generic
point of \(Y\), if and only if \(\nu(\xi_w)\) is the generic point of \(X\).

\cref{itm:ch43-id:centre} This is
\Cref{lem:ch41-support-centre-finite}\,\cref{itm:ch41-sc:centre} together
with injectivity of \(\nu\) on points.

\cref{itm:ch43-id:open} A point lies over \(U\) if and only if its scheme
point does (\Cref{lem:ch41-point-pair}), and
\(\nu^{-1}(U)^{\An}\) consists exactly of the points of \(Y^{\An}\) whose
scheme point lies in \(\nu^{-1}U\).  Restriction of a homeomorphism to
corresponding opens is a homeomorphism.

\cref{itm:ch43-id:scaling} The scaling action multiplies a semivaluation by a
positive real without changing its support, and the identification is by
restriction of absolute values; these commute, because
\((t\cdot w)|_{\kappa(\xi_v)}=t\cdot(w|_{\kappa(\xi_v)})\).  The trivial
valuation of \(K(Y)\) restricts to the trivial valuation of \(K(X)\), and by
uniqueness it is the only preimage.

\cref{itm:ch43-id:trivial} A point is \(v_{W,\triv}\) if and only if its
absolute value is trivial and its scheme point is the generic point of
\(W\), by \Cref{lem:ch41-point-pair}; both conditions are preserved and
reflected.
\end{proof}

\begin{warning}{A homeomorphism of analytifications is not an isomorphism of
varieties}
\label{warn:ch43-homeo-not-iso}
\Cref{lem:bj-universal-analytic-homeomorphism} says that \(X^{\An}\) and
\(Y^{\An}\) are the same topological space.  It does \emph{not} say that
\(X\) and \(Y\) are isomorphic, that \(\cO_X=\nu_*\cO_Y\), or that the two
psh classes coincide.  For the cuspidal cubic of
\Cref{ex:ch43-cusp-basic} the first fails because \(t\notin k[t^2,t^3]\); the
second fails for the same reason; and whether the third holds is the
theorem of this chapter, whose proof occupies
\Crefrange{sec:ch43-promotion}{sec:ch43-reflection} and is characteristic
dependent.

The distinction is not pedantic.  The class \(\PSH_X(\theta)\) is defined by
the existence of approximants \emph{on \(X\)}; the homeomorphism transports
functions but not the sheaves out of which the approximants are built.  A
reader who concludes ``the spaces are equal, so the psh classes are equal''
has assumed the theorem.  What the homeomorphism gives for free is only the
inclusion \(\nu^*\PSH_X(\theta)\subseteq\PSH_Y(\nu^*\theta)\) --- and even
that comes from functoriality of psh functions, not from topology --- plus
injectivity of \(\nu^*\).  Surjectivity is the whole content.
\end{warning}

\begin{corollary}\label{cor:bj-universal-usc}
For every function \(a:X^{\An}\to\mathbf R\cup\{-\infty\}\),
\[
                         \nu^*(a^*)=(\nu^*a)^*.
\tag{43.4}\label{eq:bj-universal-usc}
\]
\end{corollary}

\begin{proof}
Upper semicontinuous regularization is defined by the neighborhood formula.
A homeomorphism identifies the two neighborhood filters pointwise, so the
formula is unchanged.
The argument remains valid when the resulting infimum of local suprema is
\(-\infty\).

In detail: for \(x\in X^{\An}\) one has
\[
 a^*(x)=\inf_{V\ni x}\ \sup_{x'\in V}a(x'),
\]
the infimum over open neighbourhoods \(V\) of \(x\) in \(X^{\An}\).  Write
\(h:=\nu^{\An}\) and let \(y\in Y^{\An}\) with \(h(y)=x\).  Because \(h\) is
a homeomorphism, \(V\mapsto h^{-1}(V)\) is a bijection from the
neighbourhood filter of \(x\) onto that of \(y\), and
\(\sup_{x'\in V}a(x')=\sup_{y'\in h^{-1}V}a(h(y'))\).  Taking the infimum
over \(V\) gives
\(a^*(h(y))=(a\circ h)^*(y)\), which is
\eqref{eq:bj-universal-usc}.  Both suprema and the infimum are taken in
\(\Rbar\); no subtraction occurs, so the identity is valid without any
finiteness assumption on \(a\), and in particular at points where
\(a^*=-\infty\).
\end{proof}

\begin{corollary}
\label{cor:ch43-sup-star}
Let \((a_\alpha)_{\alpha\in A}\) be a nonempty family of functions
\(X^{\An}\to\mathbf R\cup\{-\infty\}\).  Then
\[
 \nu^*\Bigl(\bigl(\sup_\alpha a_\alpha\bigr)^*\Bigr)
 =\Bigl(\sup_\alpha\nu^*a_\alpha\Bigr)^*,
\]
and likewise for pointwise infima, decreasing limits and increasing limits.
The same holds for the pullback of an upper semicontinuous function: \(a\)
is usc if and only if \(\nu^*a\) is.
\end{corollary}

\begin{proof}
Pointwise suprema, infima and monotone limits are computed pointwise, and
\(\nu^{\An}\) is a bijection, so
\(\nu^*(\sup_\alpha a_\alpha)=\sup_\alpha\nu^*a_\alpha\) as functions on
\(Y^{\An}\).  Now apply \Cref{cor:bj-universal-usc} to the function
\(\sup_\alpha a_\alpha\).  The last assertion is the case \(a^*=a\),
combined with the fact that a homeomorphism preserves and reflects upper
semicontinuity.
\end{proof}

\begin{remark}[Where this corollary is spent]
\label{rem:ch43-sup-star-use}
\Cref{cor:ch43-sup-star} is not used in the present chapter at all: no
supremum of psh functions is formed anywhere below, in accordance with
\Cref{warn:ch43-firewall}.  It is recorded here because it is the step at
which \Cref{chap:bj-envelope-consequences} converts the envelope property
upstairs into the envelope property downstairs, in the form
\(\nu^*(\sup^*_i\varphi_i)=\sup^*_i\nu^*\varphi_i\); and because it is the
cleanest illustration of the principle that everything \emph{topological}
about \(\nu\) is free, while everything \emph{sheaf-theoretic} is not.
\end{remark}

\begin{proposition}
\label{prop:ch43-free-half}
In every characteristic, and for every \(\theta\in N^1(X)\), pullback
\[
 \nu^*:\PSH_X(\theta)\longrightarrow\PSH_Y(\nu^*\theta)
\]
is well defined and injective; it preserves the pointwise order, finite
maxima, addition of real constants, homogeneity, and the value at the
trivial valuation.  Moreover
\(\nu^*\cH^{\mathrm{gf}}_{\mathbf Q}(L)\subseteq
\cH^{\mathrm{gf}}_{\mathbf Q}(\nu^*L)\) for every
\(L\in\Pic(X)_{\mathbf Q}\).
\end{proposition}

\begin{proof}
Well-definedness is functoriality of psh functions along a surjective
morphism of integral projective varieties,
\cite[Theorem~4.7]{BJ22}; it is also immediate from
\Cref{def:bj-psh}, since pulling back a decreasing net
\(\varphi_i\in\cH^{\mathrm{gf}}_{\mathbf R}(L_i)\) gives a decreasing net
\(\nu^*\varphi_i\in\cH^{\mathrm{gf}}_{\mathbf R}(\nu^*L_i)\) with
\(c_1(\nu^*L_i)=\nu^*c_1(L_i)\to\nu^*\theta\), and generic finiteness is
preserved because \(\nu\) is dominant.  For the last clause: if
\(\varphi=\frac1m\max_j\{\log|s_j|+\lambda_j\}\) with
\(s_j\in H^0(X,mL)\), then
\(\nu^*\varphi=\frac1m\max_j\{\log|\nu^*s_j|+\lambda_j\}\) with
\(\nu^*s_j\in H^0(Y,\nu^*(mL))\), and the \(\nu^*s_j\) do not all vanish
identically because \(\nu\) is dominant and the \(s_j\) do not.

Injectivity is immediate from
\Cref{lem:bj-universal-analytic-homeomorphism}: if
\(\nu^*\varphi=\nu^*\varphi'\) then \(\varphi\) and \(\varphi'\) agree at
every point of \(X^{\An}\), because \(\nu^{\An}\) is surjective.  The
listed compatibilities hold because all of them are pointwise, and the value
at \(v_\triv\) is preserved by
\Cref{lem:ch43-identification-dictionary}\,\cref{itm:ch43-id:scaling}.
\end{proof}

\section{Promoting the numerical class}
\label{sec:ch43-promotion}

\subsection{The problem, and why it is not a technicality}
\label{sec:ch43-promotion-problem}

A decreasing GF regularization of an upstairs psh function generally has
rational classes \(M_i\) that converge to \(\nu^*\theta\), but the line
bundles \(M_i\) need not themselves descend.
Frobenius can descend coefficients only after their local frames have been
pulled back from \(X\).
The next lemma repairs the class without changing the function or its base
locus.

The difficulty is worth spelling out, because it is easy to mistake for a
matter of bookkeeping.  Let \(\psi\in\PSH_Y(\nu^*\theta)\).  By
\Cref{def:bj-psh} there is a decreasing net
\(\psi_i\in\cH^{\mathrm{gf}}_{\mathbf Q}(M_i)\) with
\(c_1(M_i)\to\nu^*\theta\).  Nothing in that definition says that
\(M_i\) is a pullback, and in general it is not: the homomorphism
\[
 \nu^*:\Pic(X)_{\mathbf Q}\longrightarrow\Pic(Y)_{\mathbf Q}
\]
is in general neither injective nor surjective for a nonnormal \(X\).  For
the cuspidal cubic \(C\) of \Cref{ex:ch43-cusp-basic} it is surjective but
massively non-injective: \(\Pic(C)\cong\mathbf G_a(k)\times\mathbf Z\) while
\(\Pic(\mathbf P^1)\cong\mathbf Z\), and the kernel of
\(\nu^*:\Pic(C)\to\Pic(\mathbf P^1)\) is the additive group
\(\mathbf G_a(k)\), which is the cokernel of \(\cO_C^\times\to\nu_*\cO_{\mathbf
P^1}^\times\) computed from the conductor square of
\Cref{lem:appA-conductor-square}.  Even in this most favourable case,
therefore, knowing the class of \(M_i\) upstairs determines a bundle
downstairs only up to a positive-dimensional ambiguity, and knowing only
that \(c_1(M_i)\) is \emph{close} to \(\nu^*\theta\) determines nothing at
all.  In higher dimension surjectivity fails as well: a nonnormal
unibranch surface can carry Cartier divisor classes on its normalization
that are not Cartier downstairs.

Now recall what the Frobenius argument of \Cref{sec:ch43-gf-descent} will
do.  Given \(s_j\in H^0(Y,mM_i)\), it will read off local coefficients
\(f_{U,j}\in B_U\) in a frame, raise them to the power \(q\), and glue.  The
gluing works because on an overlap the frames differ by a unit
\(g\in A_U^\times\) \emph{pulled back from \(X\)}, so that the \(q\)-th
powers of the coefficients differ by \(g^{-q}\), which is a transition
cocycle for a line bundle \emph{on \(X\)}.  If the frames differ by a unit
of \(B_U\) that is not in \(A_U\), the recipe produces a cocycle with values
in \(\cO_Y^\times\) and nothing descends.  In other words, the Frobenius
argument needs a trivializing cover and a system of frames \emph{pulled back
from \(X\)} --- that is, it needs the bundle itself, not merely its class,
to be a pullback.

There is no reason for \(M_i\) to be a pullback, and it is not permissible
to replace \(\psi_i\) by some other function in a pullback class, since the
net must continue to decrease to \(\psi\) pointwise.  What is permissible is
to enlarge the bundle by an amount that changes the class by an arbitrarily
small amount and changes the \emph{function} not at all.  That is the
content of \Cref{lem:bj-normalization-class-promotion}, and the mechanism is
the following elementary observation, which we isolate because it is used
twice more below.

\begin{lemma}
\label{lem:ch43-zero-fs}
Let \(Z\) be a projective \(k\)-variety, let \(N\) be a globally generated
line bundle on \(Z\), and let \(t_1,\ldots,t_N\in H^0(Z,N)\) generate \(N\)
at every point.  Then
\[
 \max_a\log|t_a|=0
 \qquad\text{identically on }Z^{\An}.
\]
\end{lemma}

\begin{proof}
Let \(v\in Z^{\An}\) and let \(\xi:=c(v)\) be its center, which exists
because \(Z\) is proper (\Cref{thm:ch32-center-compact}).  Choose a frame
\(e\) of \(N\) on an open neighbourhood \(U\) of \(\xi\) and write
\(t_a=h_ae\) with \(h_a\in\cO_Z(U)\).  By
\Cref{lem:ch34-frame-independence} the quantity \(\log|t_a|(v)\) does not
depend on the choice of \(e\), and equals \(-v(h_a)\), where \(v\) is
evaluated on the local ring at \(\xi\); this makes sense because
\(v\ge0\) on \(\cO_{Z,\xi}\) by definition of the center.

Since \(v(h_a)\ge0\) for every \(a\), we get \(\log|t_a|(v)\le0\) for every
\(a\), hence \(\max_a\log|t_a|(v)\le0\).  Conversely, global generation says
that the \(h_a\) generate the unit ideal of \(\cO_{Z,\xi}\), so some
\(h_{a_0}\) is a unit there; then \(v(h_{a_0})=0\) and
\(\log|t_{a_0}|(v)=0\).  Thus \(\max_a\log|t_a|(v)=0\).

The argument used no finiteness of \(v\) and no hypothesis on the support of
\(v\): it applies verbatim at points of proper support, where the identity
says in particular that the maximum is \emph{not} \(-\infty\), the
sections having no common zero.
\end{proof}

\begin{lemma}
\label{lem:ch43-ample-pullback}
Let \(f:W\to Z\) be a finite morphism of projective \(k\)-varieties and let
\(\gamma\in N^1(Z)_{\mathbf R}\) be ample.  Then \(f^*\gamma\) is ample.
Consequently \(f^*\) maps the ample cone into the ample cone and, being
linear and continuous, maps a neighbourhood of an ample class into the ample
cone.
\end{lemma}

\begin{proof}
For an integral class this is the standard fact that the pullback of an
ample line bundle along a finite morphism is ample, which follows from the
Nakai--Moishezon criterion or, more directly, from the fact that a finite
morphism is affine and \(f^*\cO(1)\) is relatively ample.  For a real class
write \(\gamma=\sum_\ell c_\ell\gamma_\ell\) with \(c_\ell>0\) and
\(\gamma_\ell\) ample integral classes, which is possible because the ample
cone is open and generated by integral ample classes; then
\(f^*\gamma=\sum_\ell c_\ell f^*\gamma_\ell\) is a positive combination of
ample classes, hence ample.  The final statement holds because \(f^*\) is
linear, hence continuous, and the ample cone is open.
\end{proof}

\begin{lemma}\label{lem:bj-normalization-class-promotion}
Suppose
\[
       \psi_i\in\cH^{\mathrm{gf}}_{\mathbf Q}(M_i),
       \qquad
       c_1(M_i)\longrightarrow\nu^*\theta
\tag{43.5}\label{eq:bj-class-promotion-input}
\]
on \(Y\).
After passing to a cofinal refinement, there are rational line bundles
\(L_i\) on \(X\) such that
\[
          c_1(L_i)\longrightarrow\theta
\quad\text{and}\quad
          \psi_i\in\cH^{\mathrm{gf}}_{\mathbf Q}(\nu^*L_i)
\tag{43.6}\label{eq:bj-class-promotion-output}
\]
as the same pointwise functions with the same common-zero loci.
\end{lemma}

\begin{proof}
The statement is an assertion about each \(\psi_i\) separately, together
with a statement about the classes; we prove both, and then explain in
\Cref{rem:ch43-cofinal} why a cofinal refinement is harmless.

Fix an ample rational line bundle \(A\) on \(X\) and a norm on
\(N^1(Y)_{\mathbf R}\).
After cofinal refinement choose positive rational numbers
\(\epsilon_i\to0\) such that
\[
        \lVert c_1(M_i)-\nu^*\theta\rVert=o(\epsilon_i).
\tag{43.7}\label{eq:bj-class-promotion-rate}
\]
Rational line-bundle classes are dense in \(N^1(X)_{\mathbf R}\).
Choose \(L_i\in\Pic(X)_{\mathbf Q}\) sufficiently close to
\(\theta+\epsilon_i c_1(A)\) that
\[
                         P_i:=\nu^*L_i-M_i
\tag{43.8}\label{eq:bj-class-promotion-error}
\]
is ample on \(Y\).
This is possible because \(\nu^*A\) is ample, the ample cone is open, and the
error in \eqref{eq:bj-class-promotion-rate} is of smaller order than
\(\epsilon_i\).
The choices also give \(c_1(L_i)\to\theta\).

In detail, the estimate behind the choice is the identity
\[
 \nu^*\bigl(\theta+\epsilon_ic_1(A)\bigr)-c_1(M_i)
 =\epsilon_i\,\nu^*c_1(A)
  +\bigl(\nu^*\theta-c_1(M_i)\bigr)
 =\epsilon_i\,\nu^*c_1(A)+o(\epsilon_i),
\]
by \eqref{eq:bj-class-promotion-rate}.  The class \(\nu^*c_1(A)\) is ample by
\Cref{lem:ch43-ample-pullback}, so \(\epsilon_i\nu^*c_1(A)\) lies in the
ample cone together with a ball of radius \(c\epsilon_i\) for a constant
\(c>0\) depending only on \(A\) and the chosen norm; the error term is
\(o(\epsilon_i)\), hence eventually inside that ball.  Since rational
classes are dense in \(N^1(X)_{\mathbf R}\) and \(\nu^*\) is continuous, we
may choose \(L_i\in\Pic(X)_{\mathbf Q}\) with
\(\lVert c_1(L_i)-\theta-\epsilon_ic_1(A)\rVert\) so small that
\(P_i:=\nu^*c_1(L_i)-c_1(M_i)\) still lies in the ample cone of \(Y\).
Finally \(\lVert c_1(L_i)-\theta\rVert\le\epsilon_i\lVert c_1(A)\rVert+o(1)\to0\).
We henceforth write \(P_i\) also for a rational line bundle representing that
class, chosen so that \(\nu^*L_i=M_i+P_i\) holds in \(\Pic(Y)_{\mathbf Q}\)
--- which is possible by construction, since \(P_i\) was \emph{defined} as
the difference \eqref{eq:bj-class-promotion-error} in \(\Pic(Y)_{\mathbf Q}\)
and only its class was constrained.

Choose \(r_i>0\) divisible enough that \(r_iP_i\) is a basepoint-free line
bundle.
Let \(t_{i1},\ldots,t_{iN_i}\) generate it.
At every centered semivaluation \(v\), all local coefficients of the
\(t_{ia}\) have nonnegative value and at least one is a unit.
Therefore
\[
                  \frac1{r_i}\max_a\log|t_{ia}|=0
\tag{43.9}\label{eq:bj-zero-fs-presentation}
\]
on all of \(Y^{\An}\), including proper-support points.
This is \Cref{lem:ch43-zero-fs} applied to \(Z=Y\) and \(N=r_iP_i\), divided
by \(r_i\); ampleness of \(P_i\) guarantees that some multiple is globally
generated, and every further multiple then is.

Tensor a GF presentation of \(\psi_i\) with the complete presentation
\eqref{eq:bj-zero-fs-presentation}.
The maximum is unchanged because the added function is zero.
The complete system has no common zero, so the old common-zero locus is
unchanged as well.
The new line-bundle class is \(M_i+P_i=\nu^*L_i\), which proves
\eqref{eq:bj-class-promotion-output}.

Explicitly: choose \(m_i>0\) with \(m_iM_i\) an honest line bundle and write
\[
 \psi_i=\frac1{m_i}\max_j\bigl\{\log|s_j|+\lambda_j\bigr\},
 \qquad
 s_j\in H^0(Y,m_iM_i),\quad\lambda_j\in\mathbf Q .
\]
Set
\[
 \widetilde s_{j,a}:=s_j^{\otimes r_i}\otimes t_{ia}^{\otimes m_i}
 \in H^0\bigl(Y,r_im_iM_i+m_ir_iP_i\bigr)
 =H^0\bigl(Y,r_im_i\,\nu^*L_i\bigr),
\]
and
\[
 \widetilde\psi_i
 :=\frac1{r_im_i}\max_{j,a}
   \bigl\{\log|\widetilde s_{j,a}|+r_i\lambda_j\bigr\}.
\]
Because \(\log|\cdot|\) turns tensor products into sums
(\Cref{lem:ch34-log-s-properties}\,\cref{itm:ch34-logs:product}) and because a
maximum over a product index set of a sum of two independent families splits,
\[
 \begin{aligned}
 \widetilde\psi_i
 &=\frac1{r_im_i}\Bigl[
   \max_j\bigl\{r_i\log|s_j|+r_i\lambda_j\bigr\}
   +\max_a\bigl\{m_i\log|t_{ia}|\bigr\}\Bigr]
 \\
 &=\frac1{m_i}\max_j\{\log|s_j|+\lambda_j\}
   +\frac1{r_i}\max_a\log|t_{ia}|,
 \end{aligned}
\]
which by \eqref{eq:bj-zero-fs-presentation} equals \(\psi_i+0=\psi_i\).  All
the manipulations are performed in \(\Rbar\) with no subtraction: the
splitting of the maximum is valid even when some \(\log|s_j|\) equals
\(-\infty\), because the second family is everywhere finite and the
convention \(-\infty+c=-\infty\) applies.  The weights \(r_i\lambda_j\) are
rational, and the sections \(\widetilde s_{j,a}\) do not all vanish
identically on \(Y\) because the \(s_j\) do not and the \(t_{ia}\) have no
common zero; so
\(\widetilde\psi_i\in\cH^{\mathrm{gf}}_{\mathbf Q}(\nu^*L_i)\), and its
common-zero locus
\[
 \bigcap_{j,a}\{\widetilde s_{j,a}=0\}
 =\Bigl(\bigcap_j\{s_j=0\}\Bigr)\cup\Bigl(\bigcap_a\{t_{ia}=0\}\Bigr)
 =\bigcap_j\{s_j=0\}
\]
is unchanged.  This is \eqref{eq:bj-class-promotion-output}.
\end{proof}

\begin{remark}[Why numerical pullback alone is insufficient]
\label{rem:ch43-numerical-insufficient}
The relation \(c_1(M_i)\approx\nu^*\theta\) is not a descent datum for local
frames.
The actual equality \(M_i+P_i=\nu^*L_i\) in the rational Picard group is what
allows the frame-change calculation in the Frobenius argument.
The zero presentation changes the actual bundle while preserving the
pointwise function.
\end{remark}

\begin{remark}[The cofinal refinement, and why it costs nothing]
\label{rem:ch43-cofinal}
Two refinements were made.  The first replaces the index set \(I\) by a
cofinal subset on which \(\lVert c_1(M_i)-\nu^*\theta\rVert\) is small enough
relative to a chosen sequence \(\epsilon_i\downarrow0\); the second is
implicit in choosing, for each \(i\), an \(r_i\) making \(r_iP_i\)
basepoint-free.  Neither affects the conclusion, for the following reason.
A decreasing net \((\psi_i)_{i\in I}\) and its restriction to a cofinal
subset \(I'\subseteq I\) have the same pointwise limit: for a decreasing net
the limit is the pointwise infimum, and
\(\inf_{i\in I}\psi_i=\inf_{i\in I'}\psi_i\) precisely because \(I'\) is
cofinal.  The conventions of the book require nets rather than sequences
here, since \(N^1(X)_{\mathbf R}\) need not admit a countable cofinal system
of approximations; cofinality is used, countability is not.  This is the
discipline recorded in \Cref{app:nets}.
\end{remark}

\begin{remark}[What is \emph{not} claimed]
\label{rem:ch43-promotion-not-claimed}
\Cref{lem:bj-normalization-class-promotion} does not claim that \(M_i\) is a
pullback, nor that \(M_i\) is close to a pullback, nor that \(\psi_i\) has a
presentation by pulled-back sections.  It claims only that \(\psi_i\), as a
function, admits a presentation whose ambient bundle is \(\nu^*L_i\).  The
sections \(\widetilde s_{j,a}\) still live upstairs and need not be pullbacks
of anything; making \emph{them} descend is the job of
\Cref{prop:bj-frobenius-gf-descent}, and it is there, not here, that the
characteristic enters.  The division of labour is worth remembering: this
section is characteristic free and purely positivity theoretic;
\Cref{sec:ch43-frobenius,sec:ch43-gf-descent} are characteristic \(p\) and
purely algebraic.
\end{remark}

\begin{example}[Class promotion on a cuspidal cubic, with numbers]
\label{ex:ch43-promotion-numbers}
Let \(C\subset\mathbf P^2\) be the cuspidal cubic and
\(\nu:\mathbf P^1\to C\) its normalization, as in
\Cref{ex:ch43-cusp-basic}, and let \(H:=\cO_C(1)\), so that
\(\nu^*H=\cO_{\mathbf P^1}(3)\) because \(C\) is a cubic and \(\nu\) is
birational of degree one onto its image.  Take
\(\theta:=\tfrac12c_1(H)\), so that \(\nu^*\theta=\tfrac32c_1(\cO(1))\) in
\(N^1(\mathbf P^1)_{\mathbf R}=\mathbf R\), where we identify a class with
its degree.

Suppose a regularizing net upstairs is presented with
\[
 M_i=\cO_{\mathbf P^1}(d_i),
 \qquad
 d_i=\frac32-\frac1{i},
\]
so \(c_1(M_i)\to\nu^*\theta\) as required, and the net is \emph{presented}
with bundles that are not pullbacks.  On this particular \(C\) a pullback
realizing each class does exist, because
\(\nu^*:\Pic(C)_{\mathbf Q}\to\Pic(\mathbf P^1)_{\mathbf Q}\) is surjective
(the conductor-square computation of \Cref{sec:ch43-promotion-problem});
what is not
given is a \emph{choice} of one, and the \(\mathbf G_a(k)\) kernel shows the
choice is not unique.  For a general nonnormal \(X\) even existence fails,
and that is the case the promotion lemma is designed for; the present
example exhibits its mechanics in a case small enough to compute.  Take
\(A:=H\), so
\(\nu^*c_1(A)=3\), and \(\epsilon_i:=1/\lceil\sqrt i\,\rceil\in\mathbf
Q_{>0}\) --- rationality matters, since
\Cref{lem:bj-normalization-class-promotion} is stated for rational
\(\epsilon_i\), and \(\epsilon_i\to0\) with
\(\lVert c_1(M_i)-\nu^*\theta\rVert=1/i=o(\epsilon_i)\), because
\(\lceil\sqrt i\,\rceil/i\to0\).  Choose
\[
 L_i:=\Bigl(\tfrac12+\epsilon_i\Bigr)H
 \in\Pic(C)_{\mathbf Q},
 \qquad
 \nu^*c_1(L_i)=\tfrac32+3\epsilon_i,
\]
so that
\[
 P_i=\nu^*c_1(L_i)-c_1(M_i)
    =3\epsilon_i+\frac1i>0,
\]
which is ample on \(\mathbf P^1\).  Any \(r_i\) with
\(r_i(3\epsilon_i+i^{-1})\in\mathbf Z_{\ge1}\) makes \(r_iP_i\) an honest
globally generated bundle \(\cO(n_i)\) with \(n_i\ge1\), generated by the
monomials \(u^cw^{n_i-c}\); and indeed
\(\max_c\log|u^cw^{n_i-c}|=0\) on \((\mathbf P^1)^{\An}\), since at any point
one of the two homogeneous coordinates is a unit at the center.  The
promoted presentation multiplies the original sections by \(r_i\) and tensors
with \(m_i\)-th powers of those monomials.  The function has not moved by so
much as one value; the ambient bundle has become
\(r_im_i\nu^*L_i\), which is what \Cref{sec:ch43-gf-descent} requires.

Note the shape of the numbers: the correction added to the class is
\(3\epsilon_i+i^{-1}\to0\), so \(c_1(L_i)\to\theta\); but at every finite
stage the correction is strictly positive, and it must be, since an ample
buffer cannot be taken to be zero.  This is the same ``buffer'' device as in
\Cref{lem:ch38-buffer}, used there for a different purpose.
\end{example}

\section{A uniform Frobenius exponent}
\label{sec:ch43-frobenius}

\subsection{Setup and the elementwise criterion}
\label{sec:ch43-frobenius-setup}

Assume for the rest of the main argument that
\(\operatorname{char}k=p>0\).
Choose a finite affine cover \(X=\bigcup_\alpha U_\alpha\) and write
\[
       A_\alpha=\cO_X(U_\alpha),
       \qquad
       B_\alpha=\cO_Y(\nu^{-1}U_\alpha).
\tag{43.10}\label{eq:bj-frobenius-affine-rings}
\]
Since \(\nu\) is finite and birational between integral schemes, every
\(A_\alpha\subset B_\alpha\) is an injective finite extension.

The finiteness deserves a sentence.  Because \(\nu\) is affine,
\(\nu^{-1}U_\alpha\) is affine, and because \(\nu\) is finite,
\(B_\alpha\) is a finite \(A_\alpha\)-module; in particular it is generated
as an \(A_\alpha\)-\emph{algebra} by finitely many elements, and we fix such
a system \(b_{\alpha1},\ldots,b_{\alpha n_\alpha}\) once and for all.
Injectivity of \(A_\alpha\to B_\alpha\) follows from dominance of \(\nu\)
between integral schemes: a nonzero regular function on \(U_\alpha\) is
nonzero at the generic point, whose preimage is the generic point of \(Y\).
(Injectivity is not used in the present section --- the criterion below is
about the \emph{image} of \(A_\alpha\) --- but it is used in
\Cref{sec:ch43-gf-descent} to identify a descended section uniquely, and it
is convenient to have recorded it.)

We shall use the criterion of
\Cref{prop:ch43-criteria}\,\cref{itm:ch43-crit:charp} in the following
elementwise form.

\begin{lemma}
\label{lem:ch43-elementwise-frobenius}
Assume \(\operatorname{char}k=p>0\) and let \(U\subseteq X\) be an affine
open.  For every \(b\in B_U\) there is an integer \(e(b)\ge0\) with
\(b^{p^{e(b)}}\in A_U\).  Moreover, if \(b^{p^e}\in A_U\) then
\(b^{p^{e'}}\in A_U\) for every \(e'\ge e\).
\end{lemma}

\begin{proof}
By \Cref{lem:ch43-uh-stability}\,\cref{itm:ch43-stab:bc} the restriction
\(\nu^{-1}U\to U\) is again a universal homeomorphism, and it is a morphism
of affine schemes, so
\Cref{prop:ch43-criteria}\,\cref{itm:ch43-crit:charp} applies to
\(A_U\to B_U\).  It yields, for each \(b\in B_U\), a \(p\)-power \(q\) with
\(b^q\) in the image of \(A_U\); since \(A_U\to B_U\) is injective, ``in the
image'' and ``in \(A_U\)'' agree, and we write \(b^q\in A_U\).  For the
monotonicity, if \(b^{p^e}=a\in A_U\) then
\(b^{p^{e+1}}=(b^{p^e})^p=a^p\in A_U\); induct.
\end{proof}

\begin{remark}[What the criterion does and does not give]
\label{rem:ch43-elementwise-not-enough}
\Cref{lem:ch43-elementwise-frobenius} produces an exponent \emph{for each
element separately}, and a priori \(\sup_be(b)=\infty\).  Nothing in the
Stacks Project criterion asserts otherwise, and for a general integral
universal homeomorphism the supremum genuinely is infinite: the extension
\(k[t^{p},t^{p^2},t^{p^3},\ldots]\subseteq\ldots\) is not of the type
occurring here, but see \Cref{ex:ch43-no-uniform-exponent} for a concrete
family of finite extensions along which the least exponent is unbounded.
The passage from an elementwise statement to a uniform one is exactly what
finiteness of \(B_U\) over \(A_U\), plus quasi-compactness of \(X\), is for,
and it is the content of the next lemma.
\end{remark}

\subsection{The uniform exponent}
\label{sec:ch43-frobenius-uniform}

\begin{lemma}
\label{lem:ch43-frobenius-subalgebra}
Let \(A\subseteq B\) be \(\mathbf F_p\)-algebras and let \(q=p^e\).  Then
\[
 B^{(q)}:=\{b\in B:b^q\in A\}
\]
is an \(A\)-subalgebra of \(B\).  Moreover
\(B^{(q)}\subseteq B^{(qp)}\), and if \(b\in B^{(q)}\) and \(c\in B^{(q)}\)
then \(bc\in B^{(q)}\) and \(b+c\in B^{(q)}\).
\end{lemma}

\begin{proof}
The map \(\Phi_q:B\to B\), \(b\mapsto b^q\), is a ring homomorphism, because
\(q\) is a power of \(p\) and \(B\) has characteristic \(p\): multiplicativity
is trivial and additivity is the Frobenius identity
\((x+y)^{p}=x^p+y^p\), iterated \(e\) times.  Hence
\(B^{(q)}=\Phi_q^{-1}(A)\) is the preimage of a subring under a ring
homomorphism, so a subring.  It contains \(A\), because \(a^q\in A\) for
\(a\in A\), and is stable under multiplication by \(A\); so it is an
\(A\)-subalgebra.  The inclusion \(B^{(q)}\subseteq B^{(qp)}\) is the
monotonicity of \Cref{lem:ch43-elementwise-frobenius}.
\end{proof}

\begin{warning}{This is the step that has no characteristic-zero analogue}
\label{warn:ch43-subalgebra-charp}
Everything in \Cref{lem:ch43-frobenius-subalgebra} rests on
\((x+y)^q=x^q+y^q\).  In characteristic zero the set
\(\{b:b^n\in A\}\) is closed under multiplication but not under addition:
for \(A=\mathbf Q[t^2,t^3]\subset B=\mathbf Q[t]\) one has \(t^2\in A\) and
\((1)^2\in A\), while \((t+1)^2=t^2+2t+1\notin A\).  There is therefore no
subalgebra to invoke, no finite system of generators to reduce to, and no
uniform exponent to extract; the substitute must be a genuinely different
identity, and it is the Roberts--Singh additive-power identity of
\Cref{chap:bj-subintegral-normalization}.  This single line is the reason
the present chapter is a few pages and the next one is a chapter.
\end{warning}

\begin{lemma}\label{lem:bj-uniform-frobenius}
There is one power \(q=p^e\), independent of \(\alpha\), such that
\[
                         b^q\in A_\alpha
            \qquad\text{for every }b\in B_\alpha
\tag{43.11}\label{eq:bj-uniform-frobenius}
\]
and every chart in the chosen finite cover.
\end{lemma}

\begin{proof}
For an affine universal homeomorphism in characteristic \(p\), the criterion
\cite[\href{https://stacks.math.columbia.edu/tag/0CNF}{Tag 0CNF}]
{stacks-project} says that for every \(b\in B_\alpha\) some \(p\)-power
belongs to \(A_\alpha\).
The additional condition \(p^e b\in A_\alpha\) in that general criterion is
automatic here because \(p=0\).

Choose finitely many \(A_\alpha\)-algebra generators of each finite algebra
\(B_\alpha\).
After replacing their individual exponents by their maximum, one power
\(q_\alpha\) works for every generator on \(U_\alpha\).
The Frobenius identity
\[
                    (x+y)^{q_\alpha}=x^{q_\alpha}+y^{q_\alpha}
\]
then shows that it works for every polynomial in the generators and hence
for every element of \(B_\alpha\).
Finally take a common maximum of the finitely many chart exponents.
Raising once more to a \(p\)-power preserves membership in \(A_\alpha\), so
the resulting \(q\) satisfies \eqref{eq:bj-uniform-frobenius} everywhere.

We repeat the argument with the three quantifiers displayed, since it is the
statement rather than the proof that carries the content.

\emph{Step one: one exponent per generator.}  Fix \(\alpha\).  By
\Cref{lem:ch43-elementwise-frobenius} each of the finitely many generators
\(b_{\alpha1},\ldots,b_{\alpha n_\alpha}\) of \(B_\alpha\) over
\(A_\alpha\) admits an exponent \(e(b_{\alpha\ell})\).

\emph{Step two: one exponent per chart.}  Put
\(e_\alpha:=\max_\ell e(b_{\alpha\ell})\) and \(q_\alpha:=p^{e_\alpha}\).  By
the monotonicity in \Cref{lem:ch43-elementwise-frobenius},
\(b_{\alpha\ell}\in B_\alpha^{(q_\alpha)}\) for every \(\ell\).  By
\Cref{lem:ch43-frobenius-subalgebra}, \(B_\alpha^{(q_\alpha)}\) is an
\(A_\alpha\)-subalgebra of \(B_\alpha\) containing all the generators; hence
\(B_\alpha^{(q_\alpha)}=B_\alpha\).  This is where characteristic \(p\) is
spent, and where the reduction from ``all of \(B_\alpha\)'' to ``finitely
many generators'' becomes legitimate.

\emph{Step three: one exponent for the cover.}  The cover is finite, so
\(e:=\max_\alpha e_\alpha\) is a well-defined integer; set \(q:=p^e\).  By
monotonicity again, \(B_\alpha^{(q_\alpha)}\subseteq B_\alpha^{(q)}\), so
\(B_\alpha^{(q)}=B_\alpha\) for every \(\alpha\).  This is
\eqref{eq:bj-uniform-frobenius}.  Finiteness of the cover is used exactly
once, here, and it is not removable
(\Cref{ex:ch43-no-uniform-exponent}).
\end{proof}

The exponent produced by \Cref{lem:bj-uniform-frobenius} is attached to the
chosen finite cover, and that is not yet what the descent argument needs.
The next corollary removes the restriction, and it is what will actually be
used.

\begin{corollary}
\label{cor:ch43-frobenius-all-affines}
Assume \(\operatorname{char}k=p>0\) and let \(q=p^e\) be as in
\Cref{lem:bj-uniform-frobenius}, produced from some finite affine cover
\((U_\alpha)\) of \(X\).  Then for \emph{every} affine open
\(U\subseteq X\),
\[
 b^q\in A_U
 \qquad\text{for every }b\in B_U .
\]
Consequently, more generally, for every open \(U\subseteq X\) the \(q\)-th
power map carries \(\cO_Y(\nu^{-1}U)\) into the image of
\(\cO_X(U)\); that is, \(\Phi_q(\nu_*\cO_Y)\subseteq\cO_X\) as subsheaves of
\(\nu_*\cO_Y\).
\end{corollary}

\begin{proof}
\emph{Step one: principal opens inside a chart.}  Let \(f\in A_\alpha\) and
let \(U=D(f)\subseteq U_\alpha\) be the corresponding principal open.  Since
\(\nu\) is affine, \(\nu^{-1}(D(f))=D(\nu^\#f)\) is the principal open of
\(\nu^{-1}U_\alpha\) determined by the image of \(f\), so
\[
 A_U=(A_\alpha)_f,
 \qquad
 B_U=(B_\alpha)_f .
\]
A general element of \(B_U\) is \(b/f^n\) with \(b\in B_\alpha\) and
\(n\ge0\).  Then
\[
 \Bigl(\frac b{f^n}\Bigr)^q=\frac{b^q}{f^{nq}},
\]
and \(b^q\in A_\alpha\) by \Cref{lem:bj-uniform-frobenius}, while
\(f^{nq}\in A_\alpha\) because \(f\in A_\alpha\).  Hence the \(q\)-th power
lies in \((A_\alpha)_f=A_U\).  Note that the localization is taken at an
element \emph{of the downstairs ring}; this is the whole trick, and it is
why the property localizes at all.

\emph{Step two: a basis of such opens.}  The opens of \(X\) that are
principal in some \(U_\alpha\) form a basis of the topology of \(X\).
Indeed, given \(x\in X\) and an open \(V\ni x\), choose \(\alpha\) with
\(x\in U_\alpha\); then \(V\cap U_\alpha\) is open in the affine scheme
\(U_\alpha\) and hence contains a principal open \(D(f)\ni x\) with
\(f\in A_\alpha\).

\emph{Step three: an arbitrary affine open.}  Let \(U\subseteq X\) be affine
open and let \(b\in B_U\).  By step two, \(U=\bigcup_\iota W_\iota\) with
each \(W_\iota\) principal in some \(U_{\alpha(\iota)}\); by
quasi-compactness of \(U\) the index set may be taken finite, though this is
not needed.  By step one applied to \(W_\iota\), the restriction
\(b|_{\nu^{-1}W_\iota}\) satisfies
\[
 \bigl(b|_{\nu^{-1}W_\iota}\bigr)^q=(b^q)|_{\nu^{-1}W_\iota}
 =\rho_\iota\bigl(a_\iota\bigr)
 \quad\text{for a (unique) }a_\iota\in A_{W_\iota},
\]
where \(\rho_\iota:A_{W_\iota}\to B_{W_\iota}\) is the structural map, which
is injective because \(\nu\) is dominant between integral schemes.  On an
overlap \(W_\iota\cap W_{\iota'}\), the two sections
\(a_\iota\) and \(a_{\iota'}\) restrict to elements of
\(\cO_X(W_\iota\cap W_{\iota'})\) with the same image in
\(\cO_Y(\nu^{-1}(W_\iota\cap W_{\iota'}))\), namely \(b^q\); by injectivity
of \(\cO_X\to\nu_*\cO_Y\), an injection of sheaves because \(\nu\) is
dominant and \(X,Y\) are integral, they agree.  Since \(\cO_X\) is a sheaf,
the \(a_\iota\) glue to a unique \(a\in\cO_X(U)=A_U\), and its image in
\(B_U\) agrees with \(b^q\) locally, hence globally, \(\cO_Y\) being a
sheaf too.  Thus \(b^q\in A_U\).

\emph{Step four: general opens.}  For an arbitrary open \(U\subseteq X\)
the same gluing argument applies with the \(W_\iota\) an affine cover of
\(U\); no affineness of \(U\) was used in steps two and three beyond
extracting the cover.  This gives the sheaf-level statement.
\end{proof}

\begin{remark}[Why the fixed cover is genuinely not enough]
\label{rem:ch43-fixed-cover-not-enough}
It is tempting to regard \Cref{cor:ch43-frobenius-all-affines} as a
restatement of \Cref{lem:bj-uniform-frobenius}, and to prove
\Cref{prop:bj-frobenius-gf-descent} using the cover \((U_\alpha)\) of
\Cref{lem:bj-uniform-frobenius} directly.  That is not possible, and the
reason is worth recording because it is the one place where a careless
reading produces a genuine gap.

The descent argument must read the sections \(s_j\in H^0(Y,\nu^*(mL))\) in
\emph{frames}, and a frame exists on \(U\) only if \(mL\) is trivial on
\(U\).  The cover \((U_\alpha)\) was chosen before \(L\) and \(m\) were
known --- indeed before the psh function was known --- and there is no
reason for \(mL|_{U_\alpha}\) to be trivial.  One would have to refine
\((U_\alpha)\) to a cover trivializing \(mL\), and refining changes the
rings \(A_\alpha\) and \(B_\alpha\).  Without
\Cref{cor:ch43-frobenius-all-affines} the exponent would have to be
recomputed on the refinement, and since the refinement depends on \(m\) and
\(L\), the exponent would depend on the stage \(i\) of the regularizing net.
That would not be fatal --- the descended functions would still be correct
pointwise, and since the denominator is \(q_im_i\) the class at stage \(i\)
would still be \(L_i\), so the descended family would still be a decreasing
net with the right classes --- but it would cost the single presentation
discipline that the rest of the argument uses: the descended sections at
stage \(i\) would live in \(q_im_iL_i\) with \(q_i\) varying, nothing would
prevent \(q_i\to\infty\), and no single downstairs Veronese multiple would
accommodate the whole net.

The localization step is therefore load-bearing.  It works, and it works
only in characteristic \(p\): the condition ``\(b^q\in A\)'' is stable under
localization at downstairs elements because \(\Phi_q\) is a ring
homomorphism commuting with localization, whereas the characteristic-zero
substitute --- an additive-power identity with fixed correction terms
\(c_1,\ldots,c_c\) --- must be re-derived on each localization, with new
correction terms, which is precisely why
\Cref{chap:bj-subintegral-normalization} must globalize its local data by
hand and with a single auxiliary twist.
\end{remark}

\begin{lemma}
\label{lem:ch43-minimal-exponent}
Assume \(\operatorname{char}k=p>0\).  The set
\[
 E(\nu):=\bigl\{e\ge0:\ b^{p^e}\in A_U
   \text{ for every affine open }U\subseteq X\text{ and every }b\in B_U\bigr\}
\]
is a nonempty up-set in \(\mathbf Z_{\ge0}\); it therefore has a least
element \(e_{\min}(\nu)\), and \(E(\nu)=\{e:e\ge e_{\min}(\nu)\}\).  One has
\(e_{\min}(\nu)=0\) if and only if \(\nu\) is an isomorphism.
\end{lemma}

\begin{proof}
Nonemptiness is \Cref{cor:ch43-frobenius-all-affines}; the up-set property
is the monotonicity in \Cref{lem:ch43-elementwise-frobenius}; a nonempty
up-set in \(\mathbf Z_{\ge0}\) has a least element.  If \(e_{\min}=0\) then
\(b=b^{p^0}\in A_U\) for all \(b\) and all \(U\), so
\(\cO_X\to\nu_*\cO_Y\) is surjective, hence an isomorphism, and \(\nu\) is
an isomorphism because it is affine.  Conversely an isomorphism has
\(e_{\min}=0\).
\end{proof}

\begin{remark}[What \(e_{\min}\) measures]
\label{rem:ch43-emin-meaning}
The invariant \(e_{\min}(\nu)\) is a crude but genuine measure of how far
\(X\) is from normal along \(\nu\).  It is \emph{not} the colength of the
conductor and not the degree of anything.  For a monomial curve singularity
\(A=k[t^a,t^b]\subset B=k[t]\) with \(\gcd(a,b)=1\), one has
\(B=A[t]\), so by \Cref{lem:ch43-frobenius-subalgebra}
\[
 b^{p^e}\in A\ \text{for all }b\in B
 \iff t^{p^e}\in A
 \iff p^e\in S:=\langle a,b\rangle,
\]
whereas the colength of the conductor is the number of gaps of \(S\),
namely \(\tfrac12(a-1)(b-1)\).  For \((a,b)=(2,3)\) the colength is \(1\)
and \(e_{\min}=1\), for every \(p\), since \(p\ge2\) lies in
\(S=\{0,2,3,4,\ldots\}\).  For \((a,b)=(p^2,p^2+1)\) the colength is
\(\tfrac12p^2(p^2-1)\), which is large, while \(e_{\min}=2\): the elements
of \(S\) below \(2p^2\) are \(0\), \(p^2\) and \(p^2+1\), so \(p^e\in S\)
forces \(e\ge2\), and \(e=2\) works.  Only the first \(p\)-power reachable
inside the semigroup matters.  The invariant is insensitive to the ground
field beyond its characteristic, and it is \emph{not} bounded in terms of
\(\dim X\) or of the multiplicity of the singularity:
\Cref{ex:ch43-no-uniform-exponent} exhibits curve singularities with
\(e_{\min}\) arbitrarily large.
\end{remark}

\subsection{Computing the exponent}
\label{sec:ch43-frobenius-examples}

\begin{example}[Frobenius descent on a cuspidal chart]
\label{ex:ch43-cusp-frobenius}
Let \(k\) be algebraically closed of characteristic \(p>0\), and consider
\[
                 A=k[t^2,t^3]\subset B=k[t].
\]
This is the normalization on the affine cuspidal chart of a projective
cuspidal cubic.  It is finite and subintegral, hence its map on spectra is a
universal homeomorphism.  The numerical semigroup generated by \(2\) and
\(3\) contains every integer at least two, so \(t^p\in A\).  For
\(b=\sum_i a_it^i\in B\), the Frobenius identity gives
\[
             b^p=\sum_i a_i^pt^{ip}\in A.
\]
Thus \Cref{lem:bj-uniform-frobenius} holds on this chart with \(q=p\).

For a centered semivaluation \(v\), the local GF term upstairs satisfies
\[
             \log|t|(v)=\frac1p\log|t^p|(v).
\]
The right-hand side is written entirely with a downstairs coefficient.  If
the support is the cusp, both sides are \(-\infty\); otherwise the equality
is the valuation identity \(v(t^p)=pv(t)\).  This is the one-generator
version of the coefficient descent in
\Cref{prop:bj-frobenius-gf-descent}.

Three further points can be checked by hand here, and each of them
prefigures a step of the general argument.

\emph{The exponent is minimal.}  We have \(t\notin A\), so \(e=0\) fails,
and \(e_{\min}(\nu)=1\) in the notation of
\Cref{lem:ch43-minimal-exponent}, independently of \(p\).  In particular the
uniform exponent for the projective cuspidal cubic is \(q=p\): the cover by
the two charts \(\{w\ne0\}\) and \(\{u\ne0\}\) has a singular point only in
the first, and the second chart is smooth, where \(\nu\) is an isomorphism
and \(e=0\) suffices; the maximum is \(1\).

\emph{Localization is visible.}  Take the downstairs element
\(f=t^2\in A\).  Then \(A_f=k[t^2,t^3][t^{-2}]=k[t^{\pm1}]\), because
\(t=t^3\cdot t^{-2}\); and \(B_f=k[t^{\pm1}]\) as well.  So on the principal
open \(D(t^2)\) the map \(\nu\) is already an isomorphism, and every element
has \(e=0\).  The exponent \(q=p\) still works there, trivially, which is
the content of step one of
\Cref{cor:ch43-frobenius-all-affines}: localizing at a downstairs element
can only make the condition easier.

\emph{Non-monomial elements are covered.}  The subalgebra argument is not
decorative.  Take \(b=1+t\in B\).  Then \(b\notin A\) and
\(b^2=1+2t+t^2\notin A\) when \(p\ne2\).  But
\(b^p=(1+t)^p=1+t^p\in A\), because \(1\in A\) and \(t^p\in A\).  Directly
computing \(\binom pi\equiv0\) for \(0<i<p\) is the Frobenius identity, and
the general statement that \(\{b:b^p\in A\}\) is a subalgebra is what
converts the single computation \(t^p\in A\) into a statement about all of
\(B\) --- \Cref{lem:ch43-frobenius-subalgebra}.
\end{example}

\begin{example}[A chart-dependent exponent does not globalize]
\label{ex:ch43-two-cusps}
Fix a prime \(p\) and let \(k\) be algebraically closed of characteristic
\(p\).  We construct a projective curve with two singular points whose naive
Frobenius exponents differ, and see exactly what breaks if one uses them
separately.

Let \(\pi:\mathbf P^1\to X\) be the map that collapses the two closed
subschemes supported at \(0\) and at \(\infty\) as follows.  Write
\(t\) for the affine coordinate near \(0\) and \(t'=1/t\) near \(\infty\).
Define \(X\) by gluing
\[
 U_1:=\Spec k[t^2,t^3],
 \qquad
 U_2:=\Spec k[t'^{\,p^2},t'^{\,p^2+1}],
\]
along the common open \(\Spec k[t^{\pm1}]\), on which both rings become
\(k[t^{\pm1}]\) --- for \(U_1\) because \(t=t^3t^{-2}\), and for \(U_2\)
because \(t'=t'^{\,p^2+1}t'^{-p^2}\).  The result is a projective curve
\(X\) with normalization \(\nu:\mathbf P^1\to X\), which is finite,
birational and bijective, hence a universal homeomorphism by
\Cref{prop:ch43-criteria}\,\cref{itm:ch43-crit:uh}.  It is unibranch, with
two singular points \(P_1=\nu(0)\) and \(P_2=\nu(\infty)\).

By the computation of \Cref{rem:ch43-emin-meaning}, the least exponents on
the two charts are
\[
 e_1=1
 \quad\text{on }U_1,
 \qquad
 e_2=2
 \quad\text{on }U_2,
\]
so \(q_1=p\) and \(q_2=p^2\), and the uniform exponent supplied by
\Cref{lem:bj-uniform-frobenius} is \(q=p^2=\max\{q_1,q_2\}\).

Now suppose one tried to descend a GF function chart by chart, using
\(q_1\) on \(U_1\) and \(q_2\) on \(U_2\).  Take \(L:=\cO_X(D)\) for a
Cartier divisor \(D\) supported in the smooth locus, with \(m=1\), and a
single section \(s\in H^0(\mathbf P^1,\nu^*L)\), so
\(\psi=\log|s|\).  Read \(s\) in frames \(e_1,e_2\) of \(L\) over \(U_1,U_2\)
pulled back from \(X\), obtaining coefficients \(f_1\in B_{U_1}\),
\(f_2\in B_{U_2}\) with
\[
 f_2=g^{-1}f_1
 \quad\text{on the overlap},
 \qquad
 g\in A_{U_1\cap U_2}^\times ,
\]
where \(g\) is the transition function of \(L\).  The chart-by-chart recipe
produces \(f_1^{q_1}\in A_{U_1}\) and \(f_2^{q_2}\in A_{U_2}\), and on the
overlap
\[
 f_2^{q_2}=g^{-q_2}f_1^{q_2}
 =g^{-q_2}\bigl(f_1^{q_1}\bigr)^{q_2/q_1},
\]
which is \emph{not} \(g^{-q_1}f_1^{q_1}\) unless \(q_1=q_2\).  In sheaf
language, \(f_1^{q_1}\) is a local section of \(q_1L\) and \(f_2^{q_2}\) is a
local section of \(q_2L\): they are sections of different line bundles, they
satisfy different cocycle conditions, and there is no object on \(X\) of
which they are the local expressions.  A ``function'' assembled from them by
writing \(\frac1{q_1}\log|f_1^{q_1}|\) on \(U_1\) and
\(\frac1{q_2}\log|f_2^{q_2}|\) on \(U_2\) is of course still equal to
\(\log|s|\) pointwise --- powering never changes the function --- but it is
not exhibited as a member of \(\cH^{\mathrm{gf}}_{\mathbf Q}(X,L)\), because
that requires \emph{one} denominator \(m\) and sections of \emph{one}
bundle \(mL\).

The repair is to raise both to the common exponent \(q=p^2\), which is
legitimate because \(f_1^{q_1}\in A_{U_1}\) implies
\(f_1^{q}=(f_1^{q_1})^{q/q_1}\in A_{U_1}\).  This is the content of
step three of the proof of \Cref{lem:bj-uniform-frobenius}, and the example
shows that the step is not a convenience but a necessity.  The reader may
also note that on this curve the exponent \(q=p^2\) is minimal: no smaller
\(p\)-power works on \(U_2\).
\end{example}

\begin{example}[Unbounded exponents, and where quasi-compactness enters]
\label{ex:ch43-no-uniform-exponent}
For \(n\ge1\) let
\[
 A_n:=k[t^{p^n},t^{p^n+1}]\subset B_n:=k[t],
\]
with \(k\) algebraically closed of characteristic \(p\).  Each inclusion is
finite and birational, and \(\Spec B_n\to\Spec A_n\) is bijective with
trivial residue extensions, hence a universal homeomorphism.  By the
semigroup computation of \Cref{rem:ch43-emin-meaning}, the least admissible
exponent for the \(n\)-th extension is exactly \(p^n\).

Consequently:
\begin{enumerate}
\item there is no exponent \(q\) serving all of these extensions
      simultaneously, so the uniformity of
      \Cref{lem:bj-uniform-frobenius} is genuinely a statement about a
      \emph{fixed} \(\nu\), not a universal constant;
\item the disjoint union
      \(\nu_\infty:\bigsqcup_{n\ge1}\Spec B_n\to\bigsqcup_{n\ge1}\Spec A_n\)
      is an affine, integral, radicial, surjective morphism --- a universal
      homeomorphism --- for which the conclusion of
      \Cref{cor:ch43-frobenius-all-affines} \emph{fails}: every affine open
      contained in one component admits an exponent, but no single exponent
      works on all of them.  The failure is not a failure of the argument;
      it is the absence of quasi-compactness, which is what step three of
      the proof of \Cref{lem:bj-uniform-frobenius} consumes.
\end{enumerate}

Within the standing hypotheses of this chapter neither pathology can occur:
\(X\) is a projective variety, hence quasi-compact and of finite type, and
\(\nu\) is finite.  Moreover, by
\Cref{prop:ch43-criteria}\,\cref{itm:ch43-crit:finite} the finiteness is not
an extra hypothesis: a universal homeomorphism of finite type is
automatically finite.  What one may \emph{not} do is bound \(e_{\min}(\nu)\)
in terms of \(p\), \(\dim X\), the degree of \(X\) in a fixed embedding, or
any other discrete invariant fixed in advance --- the curves above may be
projectivized inside \(\mathbf P^{p^n+1}\) and have \(e_{\min}=n\).
\end{example}

\begin{remark}[The characteristic-zero contrast, stated once more]
\label{rem:ch43-charzero-contrast}
Over a field of characteristic zero the naive analogue of
\Cref{cor:ch43-frobenius-all-affines} is false already in the simplest
case.  Take \(A=\mathbf Q[t^2,t^3]\subset B=\mathbf Q[t]\), a finite
subintegral extension, and ask for an integer \(q\ge1\) with \(b^q\in A\)
for every \(b\in B\).  The element \(b=t\) passes the test for every
\(q\ge2\), which is why the difficulty is invisible if one tests only
monomials; but \(b=1+t\) fails for every \(q\), since
\[
 (1+t)^q=1+qt+\binom q2t^2+\cdots
\]
has the term \(qt\) with \(q\ne0\) in \(\mathbf Q\), and no other term
contributes a multiple of \(t\); hence \((1+t)^q\notin A\) for every
\(q\ge1\).  The obstruction is exactly the failure of
\Cref{lem:ch43-frobenius-subalgebra}: in characteristic zero the set
\(\{b:b^q\in A\}\) is not closed under addition, so it is not a subalgebra,
so the reduction to generators is unavailable, so there is no uniform
exponent to be had.

The correct characteristic-zero replacement, due to Roberts and Singh, is to
allow the power to be corrected by lower-order terms with \emph{fixed}
coefficients: with \(A\subseteq B\) subintegral there are
\(c_1,\ldots,c_c\in B\) and \(N\) such that
\(b^n+\sum_{i}\binom nic_ib^{n-i}\in A\) for all \(n\ge N\).  For
\(A=\mathbf Q[t^2,t^3]\) and \(b=1+t\), one may take \(c=1\) and
\(c_1=-t\), since
\[
 (1+t)^n-n\,t\,(1+t)^{n-1}
 =(1+t)^{n-1}\bigl(1+t-nt\bigr),
\]
whose expansion has vanishing coefficient in degree one:
\((n-1)+(1-n)=0\).  So the corrected power lies in \(A\) for every
\(n\ge1\).  Everything downstream --- systems instead of single sections, a
valuative sandwich instead of an identity, a limit \(n\to\infty\) instead of
an equality, an artificial monotonization --- follows from the presence of
those correction terms.  See \Cref{chap:bj-subintegral-normalization}, and
\Cref{thm:appA-roberts-singh} for the statement of the theorem itself.
\end{remark}

\section{Descent of a GF stage}
\label{sec:ch43-gf-descent}

We come to the step for which the previous two sections were preparation.
The statement is the exact reflection statement for Fubini--Study functions:
a rational GF function upstairs, in a class pulled back from \(X\), is
already a rational GF function downstairs.  No approximation, no error term,
no passage to a limit; an identity of functions on the nose.

Before the proof, it is worth naming the three things that could go wrong
and observing that each is handled by one of the preparations.  The
coefficients could fail to lie downstairs --- handled by
\Cref{cor:ch43-frobenius-all-affines}.  The powered coefficients could fail
to glue --- handled by the observation that raising a cocycle to the power
\(q\) gives a cocycle, of the \(q\)-th Veronese multiple.  The identity could
fail at the common zeros of the sections, where both sides are \(-\infty\)
--- handled by the fact that powering a section neither creates nor destroys
zeros, so both sides are \(-\infty\) at exactly the same points.

\begin{proposition}\label{prop:bj-frobenius-gf-descent}
Let \(L\in\Pic(X)_{\mathbf Q}\), and assume
\(\operatorname{char}k=p>0\).
Every rational GF function
\[
                         \psi\in
              \cH^{\mathrm{gf}}_{\mathbf Q}(\nu^*L)
\]
is the pullback of a rational GF function in the exact class \(L\).
The equality holds on all analytic points, including the common
\(-\infty\)-locus.
\end{proposition}

\begin{proof}
Choose \(m>0\) sufficiently divisible that \(mL\) is a line bundle and write
\[
 \psi=\frac1m\max_j\{\log|s_j|+\lambda_j\},
 \qquad
 s_j\in H^0(Y,\nu^*(mL)),
 \quad
 \lambda_j\in\mathbf Q.
\tag{43.12}\label{eq:bj-positive-gf}
\]
Fix the uniform Frobenius power \(q\) from
\Cref{lem:bj-uniform-frobenius}.
Choose a finite cover of \(X\) by affine opens \(U_\alpha\) \emph{on which
\(mL\) is trivial}; this cover depends on \(L\) and on \(m\) and need have
nothing to do with the cover used to produce \(q\).  The exponent \(q\)
nevertheless applies to it, by
\Cref{cor:ch43-frobenius-all-affines}: that corollary is exactly what makes
the present sentence legitimate, and without it the exponent would have to
be recomputed here, with the consequences described in
\Cref{rem:ch43-fixed-cover-not-enough}.
On \(U_\alpha\), trivialize \(mL\) by a frame \(e_\alpha\) and write
\[
                         s_j=f_{\alpha j}\nu^*e_\alpha,
             \qquad f_{\alpha j}\in B_\alpha.
\]
By \eqref{eq:bj-uniform-frobenius},
\(f_{\alpha j}^q\in A_\alpha\).

Here \(\nu^*e_\alpha\) is a frame of \(\nu^*(mL)\) over \(\nu^{-1}U_\alpha\)
--- pullback of a frame is a frame, since \(\nu^*\) is a monoidal functor on
line bundles --- and \(f_{\alpha j}\) is the unique element of
\(B_\alpha=\cO_Y(\nu^{-1}U_\alpha)\) with \(s_j=f_{\alpha j}\nu^*e_\alpha\).
It is here, and only here, that the conclusion of
\Cref{lem:bj-normalization-class-promotion} is used: had the ambient bundle
of \(\psi\) been an arbitrary \(M\) with \(c_1(M)\) merely close to
\(\nu^*c_1(L)\), there would be no pulled-back frame \(\nu^*e_\alpha\) to
read \(s_j\) against, and the following cocycle computation would be
meaningless.

If \(e_\beta=g_{\alpha\beta}e_\alpha\) on an overlap, then
\[
 f_{\beta j}=g_{\alpha\beta}^{-1}f_{\alpha j},
 \qquad
 f_{\beta j}^q=g_{\alpha\beta}^{-q}f_{\alpha j}^q.
\]
These are precisely the transition laws for sections of \(mqL\).
The local \(q\)-th powers therefore glue to sections
\[
          s_j^{[q]}\in H^0(X,mqL),
          \qquad
          \nu^*s_j^{[q]}=s_j^q.
\tag{43.13}\label{eq:bj-frobenius-sections}
\]
Define
\[
 \varphi:=\frac1{mq}\max_j
       \{\log|s_j^{[q]}|+q\lambda_j\}
       \in\cH^{\mathrm{gf}}_{\mathbf Q}(L).
\tag{43.14}\label{eq:bj-positive-descended-gf}
\]
For every analytic point,
\[
 \nu^*\varphi
 =\frac1{mq}\max_j
       \{\log|s_j^q|+q\lambda_j\}
 =\frac1m\max_j\{\log|s_j|+\lambda_j\}
 =\psi.
\tag{43.15}\label{eq:bj-positive-gf-identity}
\]
If all \(s_j\) vanish at the point, every term in both maxima is
\(-\infty\).
The identity is therefore an equality in
\(\mathbf R\cup\{-\infty\}\), not merely an equality away from the base
locus.

Five verifications remain, and we record them explicitly.

\emph{The gluing.}  On \(U_\alpha\cap U_\beta\) the elements
\(f_{\alpha j}^q\) and \(f_{\beta j}^q\) lie in the corresponding rings
\(A_{U_\alpha}\), \(A_{U_\beta}\) by
\Cref{cor:ch43-frobenius-all-affines}, and satisfy
\(f_{\beta j}^q=g_{\alpha\beta}^{-q}f_{\alpha j}^q\) with
\(g_{\alpha\beta}\in\cO_X^\times(U_\alpha\cap U_\beta)\), because
\(g_{\alpha\beta}\) is the transition function of the bundle \(mL\)
\emph{on \(X\)}.  Since \((g_{\alpha\beta}^{q})\) is the transition cocycle
of \(qmL\), the collection \((f_{\alpha j}^q e_\alpha^{\otimes q})_\alpha\)
is a compatible family of local sections of \(qmL\) and glues to a global
section, which we call \(s_j^{[q]}\).  Uniqueness of the glued section
follows from the sheaf axiom.

\emph{The compatibility \(\nu^*s_j^{[q]}=s_j^{\otimes q}\).}  On
\(\nu^{-1}U_\alpha\) both sides equal
\(f_{\alpha j}^q\,(\nu^*e_\alpha)^{\otimes q}\), and a global section of a
sheaf is determined by its restrictions.

\emph{The classes.}  \(s_j^{[q]}\in H^0(X,qmL)\), so
\(\frac1{qm}\max_j\{\log|s_j^{[q]}|+q\lambda_j\}\) is a rational GF function
for \(L\), the weights \(q\lambda_j\) being rational.

\emph{Generic finiteness.}  The \(s_j\) do not all vanish identically on
\(Y\), so some \(f_{\alpha j}\ne0\) for some \(\alpha\), so
\(f_{\alpha j}^q\ne0\) --- \(B_\alpha\) being a domain --- and hence
\(s_j^{[q]}\ne0\).  Thus \(\varphi\) is not identically \(-\infty\), and it
belongs to \(\cH^{\mathrm{gf}}_{\mathbf Q}(L)\) in the sense of
\Cref{def:gf-fs-functions}.

\emph{The pointwise identity, in \(\Rbar\).}  Fix \(w\in Y^{\An}\) and let
\(v=\nu^{\An}(w)\).  By
\Cref{lem:ch41-support-centre-finite}\,\cref{itm:ch41-sc:sections} the value
\(\log|s_j^{[q]}|(v)\) may be computed in a frame near \(c(v)\), and by
\Cref{lem:ch43-identification-dictionary}\,\cref{itm:ch43-id:centre} we have
\(c(w)=\nu^{-1}(c(v))\), so the frames \(e_\alpha\) and \(\nu^*e_\alpha\)
match.  Then
\[
 \log\bigl|s_j^{[q]}\bigr|(v)
 =-v\bigl(f_{\alpha j}^q\bigr)
 =-w\bigl(f_{\alpha j}^q\bigr)
 =-q\,w(f_{\alpha j})
 =q\log|s_j|(w),
\]
where the second equality is the identification of \(v\) with \(w\) and the
third is multiplicativity of the semivaluation.  All four quantities lie in
\([-\infty,+\infty)\), and the identity \(w(f^q)=q\,w(f)\) holds also when
\(w(f)=+\infty\), that is, when \(f\) lies in the support of \(w\), since
\(q\cdot(+\infty)=+\infty\).  Dividing by \(qm\) and adding \(\lambda_j\)
gives \eqref{eq:bj-positive-gf-identity} term by term, and the maximum of
equal families is equal.  No subtraction of infinite values occurs at any
stage.
\end{proof}

\begin{remark}[Why weights are also multiplied by \(q\)]
\label{rem:ch43-weights-times-q}
The equality \(\log|s_j^q|=q\log|s_j|\) accounts for the section term.
Multiplying \(\lambda_j\) by the same \(q\) makes the whole maximum scale by
\(q\), which is then cancelled by the denominator \(mq\).
Leaving the weights unchanged would alter the function.

Concretely, with unchanged weights one would obtain
\(\frac1{qm}\max_j\{q\log|s_j|+\lambda_j\}\), which differs from \(\psi\)
unless all the \(\lambda_j\) are equal; and when they are not, the
difference is not even a constant, since the maximum is attained by
different indices at different points.  The rescaling
\(\lambda_j\mapsto q\lambda_j\) is the unique one making the presentation
homogeneous of degree \(q\) in the pair (section, weight), which is what the
denominator \(qm\) then undoes.
\end{remark}

\begin{remark}[No averaging]
\label{rem:ch43-no-averaging}
The construction takes powers of coefficients; it does not average over
fibres and does not divide in the ground field.
It therefore remains valid for a purely inseparable map.
The only division by \(q\) occurs in the real normalization of the GF
function.

The contrast with \Cref{chap:bj-finite-descent} is instructive.  There, a
finite morphism \(h:T\to X\) of degree \(d\) descends a GF function by
forming the norm of a universal linear form and dividing by \(md\); the
resulting function is the \emph{fibre average}
\(\frac1d\sum_{w\mid z}e(w/z)\varphi(w)\), which in general is not equal to
\(\varphi\) at any lift.  Here every fibre is a single point with
multiplicity one in the relevant sense, the ``average'' would be the value
itself, and the mechanism is completely different: it produces the
\emph{same} function, not an averaged one.  In particular
\Cref{prop:bj-frobenius-gf-descent} is not a special case of
\Cref{prop:bj-weighted-coefficient-norm}, and the two are proved by
unrelated means.  The number \(q\) is an exponent, not a degree; there is no
division by \(\deg\nu=1\) anywhere, and no separability hypothesis is
required or available.
\end{remark}

\begin{corollary}
\label{cor:ch43-gf-bijection}
Assume \(\operatorname{char}k=p>0\) and let
\(L\in\Pic(X)_{\mathbf Q}\).  Then
\[
 \nu^*:\cH^{\mathrm{gf}}_{\mathbf Q}(X,L)
 \longrightarrow
 \cH^{\mathrm{gf}}_{\mathbf Q}(Y,\nu^*L)
\]
is a bijection.  It is an isomorphism for the pointwise order, commutes with
finite maxima and with addition of rational constants, and matches
\(-\infty\)-loci: for \(\varphi\) downstairs with common-zero locus
\(Z\subseteq X\), the pullback \(\nu^*\varphi\) has common-zero locus
\(\nu^{-1}Z\).
\end{corollary}

\begin{proof}
Injectivity and well-definedness are
\Cref{prop:ch43-free-half}; surjectivity is
\Cref{prop:bj-frobenius-gf-descent}.  The order and finite-maximum statements hold
because both maps are given by composition with the bijection
\(\nu^{\An}\).  For the loci, \(\varphi=-\infty\) exactly on the analytic
points whose support lies in the common zero locus of a presentation
(\Cref{prop:ch34-generic-finiteness}\,\cref{itm:ch34-gf:polar}), and supports
correspond under
\Cref{lem:ch43-identification-dictionary}\,\cref{itm:ch43-id:supp}.
\end{proof}

\begin{remark}[Bijective on functions, not on presentations]
\label{rem:ch43-not-presentations}
\Cref{cor:ch43-gf-bijection} is a bijection between \emph{functions}.  It is
emphatically not a bijection between presentations, nor an isomorphism
\(H^0(X,mL)\to H^0(Y,\nu^*(mL))\) --- the latter is injective but usually
far from surjective.  For the cuspidal cubic \(C\) with \(L=\cO_C(1)\) and
\(m=1\), the space \(H^0(C,\cO(1))\) is three-dimensional while
\(H^0(\mathbf P^1,\cO(3))\) is four-dimensional, and the missing section is
precisely the one whose local coefficient at the cusp is \(t\).  What
\Cref{prop:bj-frobenius-gf-descent} provides is not that section but a
\emph{different} presentation, in the Veronese multiple \(qmL\), of the same
function.  The denominator \(m\) is not preserved; only the class \(L\) and
the function are.
\end{remark}

\begin{proposition}
\label{prop:ch43-frobenius-no-birational}
Let \(\operatorname{char}k=p>0\) and let \(\mu:Y\to X\) be \emph{any} finite
universal homeomorphism of integral projective \(k\)-varieties, not assumed
birational.  Then \Cref{lem:bj-uniform-frobenius},
\Cref{cor:ch43-frobenius-all-affines} and
\Cref{prop:bj-frobenius-gf-descent} hold verbatim for \(\mu\), and so does
\Cref{lem:bj-normalization-class-promotion}.  Consequently
\Cref{thm:bj-positive-universal-reflection} holds for \(\mu\).
\end{proposition}

\begin{proof}
We inspect the four proofs.

\Cref{lem:bj-normalization-class-promotion} uses only that \(\mu\) is finite,
through \Cref{lem:ch43-ample-pullback}, and that \(Y\) is projective,
through \Cref{lem:ch43-zero-fs}.

\Cref{lem:bj-uniform-frobenius} uses
\Cref{prop:ch43-criteria}\,\cref{itm:ch43-crit:charp}, which requires only
that \(\mu\) be an affine universal homeomorphism of
\(\mathbf F_p\)-schemes, and finiteness of \(\mu\), to have finitely many
algebra generators.  Injectivity of \(A_U\to B_U\) was used only to say
``\(b^q\in A_U\)'' instead of ``\(b^q\) lies in the image of \(A_U\)''; for
\(X\) and \(Y\) integral and \(\mu\) dominant --- which holds since \(\mu\)
is surjective --- injectivity is automatic and independent of
birationality.

\Cref{cor:ch43-frobenius-all-affines} uses the same, plus the sheaf axiom
and injectivity of \(\cO_X\to\mu_*\cO_Y\), again a consequence of dominance
between integral schemes.

\Cref{prop:bj-frobenius-gf-descent} uses the previous two, the cocycle
computation --- which involves only the transition functions of a bundle on
\(X\) --- and \Cref{lem:bj-universal-analytic-homeomorphism}, whose proof
used radiciality and surjectivity but not birationality.

Finally the proof of \Cref{thm:bj-positive-universal-reflection} assembles
these with \Cref{def:bj-psh} and
\Cref{thm:bj-basic-closure}\,(5), neither of which mentions
birationality.
\end{proof}

\begin{remark}[Why the hypothesis is nevertheless stated]
\label{rem:ch43-birational-still-stated}
Three reasons.  First, in characteristic zero birationality is automatic: a
finite universal homeomorphism in characteristic zero has trivial residue
extensions (\Cref{rem:appA-subintegral-stronger}), in particular at the
generic point, so \(K(Y)=K(X)\).  Stating \eqref{eq:bj-universal-map} with
birationality therefore costs nothing there and makes the two halves of
\Cref{thm:bj-positive-universal-reflection} and its characteristic-zero
counterpart symmetric.  Second, the application is to a normalization, which
is birational by construction.  Third,
\Cref{prop:ch43-frobenius-no-birational} is a genuine strengthening only in
characteristic \(p\), and it is worth having isolated:
\Cref{ex:ch43-frobenius-uh} shows the extra generality is not empty, and
\Cref{ex:ch43-frobenius-reflection} carries it out on \(\mathbf P^1\).
\end{remark}

\subsection{The descent, computed}
\label{sec:ch43-gf-descent-example}

\begin{example}[Descending a Fubini--Study function through the cusp]
\label{ex:ch43-cusp-gf-descent}
Let \(k\) be algebraically closed of characteristic \(p>0\), let
\(C\subset\mathbf P^2\) be the cuspidal cubic
\(\{ZY^2=X^3\}\) and \(\nu:\mathbf P^1\to C\) its normalization, as in
\Cref{ex:ch43-cusp-basic}.  Cover \(C\) by the two affine opens
\[
 U_1:=C\cap\{Z\ne0\},
 \qquad
 U_2:=C\cap\{Y\ne0\},
\]
which do cover: a point of \(C\) with \(Y=Z=0\) would satisfy \(X^3=0\).
Write \(t=u/w\) for the affine coordinate on \(\nu^{-1}U_1=\{w\ne0\}\) and
\(t'=1/t\) on \(\nu^{-1}U_2=\{u\ne0\}\).  Then
\[
 A_{U_1}=k[t^2,t^3],\quad B_{U_1}=k[t],
 \qquad
 A_{U_2}=k[t',t'^{\,3}]=k[t']=B_{U_2},
\]
the second because \(X/Y=t'\) and \(Z/Y=t'^{\,3}\), so that \(\nu\) is an
isomorphism over \(U_2\).  By \Cref{ex:ch43-cusp-frobenius} the uniform
exponent is \(q=p\).

Take \(L:=\cO_C(1)\), so \(\nu^*L=\cO_{\mathbf P^1}(3)\) --- because \(C\)
has degree three and \(\nu\) is birational onto its image --- and \(m=1\).
Let \(e_1,e_2\) be the frames of \(\cO_C(1)\) over \(U_1,U_2\) attached to
the coordinates \(Z\) and \(Y\); the coefficient of a section \(s\) is then
\(s/Z\) on \(U_1\) and \(s/Y\) on \(U_2\).

The space \(H^0(\mathbf P^1,\cO(3))\) has basis \(u^3,u^2w,uw^2,w^3\), and
the pullbacks of the coordinate sections are
\[
 \nu^*X=u^2w,
 \qquad
 \nu^*Y=u^3,
 \qquad
 \nu^*Z=w^3 .
\]
The fourth basis vector \(uw^2\) is not a pullback: its coefficient in the
frame \(\nu^*e_1\) is
\[
 \frac{uw^2}{w^3}=t\notin A_{U_1},
\]
whereas the three pullbacks have coefficients \(t^2,t^3,1\), all in
\(A_{U_1}\).  So take the upstairs GF function
\[
 \psi:=\max\bigl\{\log|uw^2|,\ \log|u^3|+\lambda\bigr\},
 \qquad
 \lambda\in\mathbf Q,
\]
whose coefficients on \(\nu^{-1}U_1\) are
\(f_{1,1}=t\) and \(f_{1,2}=t^3\).  Since \(f_{1,1}\notin A_{U_1}\), the
function \(\psi\) is visibly a GF function upstairs and not visibly one
downstairs.

Apply the recipe.  Raise to the power \(q=p\):
\[
 f_{1,1}^p=t^p\in A_{U_1},
 \qquad
 f_{1,2}^p=t^{3p}\in A_{U_1},
\]
using \(t^p=(t^2)^a(t^3)^b\) for any \(a,b\ge0\) with \(2a+3b=p\), which
exist because \(p\ge2\); on \(U_2\) the coefficients are
\(uw^2/u^3=t'^{\,2}\) and \(u^3/u^3=1\), whose \(p\)-th powers lie in
\(A_{U_2}=k[t']\) trivially.  The glued sections are
\[
 s_1^{[p]}=X^aY^bZ^{\,p-a-b}\big|_C,
 \qquad
 s_2^{[p]}=Y^p\big|_C,
\]
both in \(H^0(C,\cO_C(p))\); their coefficients in the frame
\(e_1^{\otimes p}\) are \(x^ay^b=t^{2a+3b}=t^p\) and \(y^p=t^{3p}\), as
required, and one checks \(a+b\le p\) since \(2(a+b)\le2a+3b=p\).  The
descended function is
\[
 \varphi=\frac1p\max\bigl\{\log\bigl|s_1^{[p]}\bigr|,\
   \log\bigl|s_2^{[p]}\bigr|+p\lambda\bigr\}
 \in\cH^{\mathrm{gf}}_{\mathbf Q}\bigl(C,\cO_C(1)\bigr),
\]
and \(\nu^*\varphi=\psi\).

Let us check three points explicitly.  Let \(v_a\) denote the monomial
semivaluation on \(\mathbf P^1\) with \(v_a(t)=a\ge0\) and
\(v_a(w/w)=0\), normalized on the chart \(U_1\); it is a full-support point
for \(a<\infty\), and \(a=\infty\) corresponds to the point of proper
support at the cusp preimage.

\emph{At a generic point} (\(a=0\), the trivial valuation of
\(k(t)\)): \(\psi(v_0)=\max\{0,\lambda\}\) and
\(\varphi(v_0)=\frac1p\max\{0,p\lambda\}=\max\{0,\lambda\}\).  Equal.

\emph{At a divisorial point} (\(a=1\), that is, \(\ord_{t=0}\)):
\(\log|uw^2|(v_1)=-1\) and \(\log|u^3|(v_1)=-3\), so
\(\psi(v_1)=\max\{-1,-3+\lambda\}\).  Downstairs
\(\log|s_1^{[p]}|(v_1)=-p\) and \(\log|s_2^{[p]}|(v_1)=-3p\), so
\(\varphi(v_1)=\frac1p\max\{-p,-3p+p\lambda\}=\max\{-1,-3+\lambda\}\).
Equal.  For \(a>0\) general the same computation gives
\(\max\{-a,-3a+\lambda\}\) on both sides.

\emph{At the point of proper support at the cusp} (\(a=\infty\), i.e. the
semivaluation whose support is the maximal ideal at \(t=0\)): both sections
vanish there, so \(\psi=-\infty\); and both descended sections vanish at the
cusp of \(C\), so \(\varphi=-\infty\).  Equal, in \(\Rbar\).  This is the
point that a proof working only on divisorial valuations would have
discarded, and it is retained here for free, because powering a section does
not move its zero locus.
\end{example}

\begin{example}[Reflection through a Frobenius, with numbers]
\label{ex:ch43-frobenius-reflection}
This example illustrates
\Cref{prop:ch43-frobenius-no-birational}: a finite universal homeomorphism
that is not birational, with an explicit Fubini--Study function carried
through the reflection in both directions.

Let \(k\) be algebraically closed of characteristic \(p>0\), let
\(X=\mathbf P^1\) with coordinate \(t\), and let \(Y=\mathbf P^1\) with
coordinate \(\tau\), and let
\[
 \mu:Y\longrightarrow X,
 \qquad
 \mu^\#(t)=\tau^p .
\]
By \Cref{ex:ch43-frobenius-uh} this is a finite universal homeomorphism of
degree \(p\), not birational and not an isomorphism.  Put
\(L:=\cO_X(1)\), so \(\mu^*L=\cO_Y(p)\).

\emph{The exponent.}  On the chart \(U=\{w\ne0\}\) with
\(A_U=k[t]\hookrightarrow B_U=k[\tau]\), \(t\mapsto\tau^p\), the algebra
\(B_U\) is generated over \(A_U\) by \(\tau\), and \(\tau^p=t\in A_U\), so
\(q=p\) works and \(e_{\min}(\mu)=1\).  The same on the other chart.

\emph{Downstairs to upstairs.}  Take
\(\varphi:=\max\{\log|t|,\,\lambda\}\), a GF function for \(L\) in the
presentation with sections \(u,w\) of \(\cO_X(1)\) and weights \(0,\lambda\)
--- writing \(t=u/w\).  Then
\[
 \mu^*\varphi=\max\{\log|\tau^p|,\lambda\}
             =\max\{p\log|\tau|,\lambda\},
\]
which is a GF function for \(\mu^*L=\cO_Y(p)\) presented by
\(u_Y^p,w_Y^p\) with weights \(0,\lambda\).

\emph{Upstairs to downstairs.}  Take
\(\psi:=\max\{\log|u_Y^p|,\ \log|w_Y^p|+\lambda'\}\), a GF function for
\(\mu^*L\).  Its coefficient in the frame \(\mu^*e\) on the chart is
\(f=\tau^p\), and \(f^{q}=\tau^{p^2}=t^p\in A_U\).  The glued section is
\(s^{[p]}\) with local expression \(t^p\), that is, \(u^{p}\) as a section of
\(\cO_X(p)=pL\); similarly for the second.  The descended function is
\[
 \frac1p\max\bigl\{\log|u^p|,\ \log|w^p|+p\lambda'\bigr\}
 =\max\bigl\{\log|u|,\ \log|w|+\lambda'\bigr\},
\]
a GF function for \(L\) whose pullback is \(\psi\).

\emph{The upshot for classes.}  \(\mu^*\) carries
\(\PSH_X(\theta)\) bijectively onto \(\PSH_Y(\mu^*\theta)=\PSH_Y(p\theta)\)
by \Cref{thm:bj-positive-universal-reflection} together with
\Cref{prop:ch43-frobenius-no-birational}.  Under the identification
\(X^{\An}=Y^{\An}\) of
\Cref{lem:bj-universal-analytic-homeomorphism}, \(\mu^{\An}\) is the scaling
of \Cref{prop:ch32-scaling-action} by the factor \(p\): a semivaluation
\(w\) on \(k(\tau)\) restricts to the semivaluation \(t\mapsto p\,w(\tau)\)
on \(k(t)\).  So the bijection here is not the identity map on functions, as
it is in the birational case, but a reparametrization of the cone
\(X^{\An}\) by a factor \(p\), together with a rescaling of the class by
\(p\).  The two effects cancel, which is what the reflection theorem
asserts.
\end{example}

\section{Psh reflection in positive characteristic}
\label{sec:ch43-reflection}

The word \emph{reflection} is used here in the sense fixed in
\Cref{chap:bj-reflection}: a statement that pullback along a morphism is a
\emph{bijection} of psh classes, so that the two classes may be identified
and a property proved upstairs may be read downstairs.  The reflection
theorems of the book differ in what they require and in what they deliver,
and it is worth placing the present one among them before proving it.

\Cref{thm:bj-homogeneous-reflection} reflects \emph{homogeneous} psh
functions along a projective birational morphism of normal varieties, and
\Cref{thm:bj-carrier-reflection} extends this to psh functions dominated by
a homogeneous carrier; both are proved by realizing b-divisors and both
require normality on both sides.  \Cref{cor:ch41-finite-bijection} reflects
along a finite morphism of normal varieties, but only onto the fibrewise
constant functions, and it is a consequence of the envelope property rather
than an input to it.  The present theorem is of a third kind: no
normality downstairs, no carrier hypothesis, no envelope property, and a
bijection onto the whole class.  What buys all of that is that
\(\nu^{\An}\) is a homeomorphism, so that ``the same function'' is a
meaningful phrase, and that in characteristic \(p\) a presentation can be
converted rather than approximated.

\begin{theorem}\label{thm:bj-positive-universal-reflection}
Assume \(\operatorname{char}k=p>0\).
For every \(\theta\in N^1(X)_{\mathbf R}\), pullback is a bijection
\[
 \nu^*:\PSH_X(\theta)
       \xrightarrow{\ \sim\ }
       \PSH_Y(\nu^*\theta).
\tag{43.16}\label{eq:bj-positive-universal-reflection}
\]
\end{theorem}

\begin{proof}
Dominant pullback preserves psh functions and numerical classes by
\cite[Theorem~4.7]{BJ22}.
It is injective because \(\nu^{\An}\) is surjective, indeed a homeomorphism.

For surjectivity, let
\[
                         \psi\in\PSH_Y(\nu^*\theta).
\]
Choose a decreasing rational GF regularization
\[
 \psi_i\in\cH^{\mathrm{gf}}_{\mathbf Q}(M_i),
 \qquad
 \psi_i\downarrow\psi,
 \qquad
 c_1(M_i)\longrightarrow\nu^*\theta.
\tag{43.17}\label{eq:bj-positive-upstairs-regularization}
\]
Apply \Cref{lem:bj-normalization-class-promotion}.
After cofinal refinement, the same functions have exact classes
\(M_i=\nu^*L_i\), where \(c_1(L_i)\to\theta\).

By \Cref{prop:bj-frobenius-gf-descent}, write
\(\psi_i=\nu^*\varphi_i\) with
\[
                         \varphi_i\in
                         \cH^{\mathrm{gf}}_{\mathbf Q}(L_i).
\]
Under the analytic homeomorphism, the \(\varphi_i\) are the same pointwise
functions as the \(\psi_i\).
They therefore decrease to the function corresponding to \(\psi\), at every
analytic point and on every common-zero locus.
Their classes converge to \(\theta\).

Apply the decreasing-limit axiom PSH2.
Its exceptional alternative is an identically \(-\infty\) limit.
That alternative cannot occur because \(\psi\) is psh and hence has a finite
value at the trivial valuation; equivalently, each homogeneous normalization
used in the regularization has the prescribed finite trivial value.
Thus the pointwise limit belongs to \(\PSH_X(\theta)\).
Its pullback is \(\psi\), proving surjectivity.

We restate the last three paragraphs with all the references discharged, in
order to make the audit of \Cref{sec:ch43-audit} possible by inspection.

\emph{The regularization exists.}  By \Cref{thm:ch34-fs-regularization}
applied on \(Y\), the function \(\psi\in\PSH_Y(\nu^*\theta)\) is the
pointwise limit of a decreasing net
\(\psi_i\in\cH^{\mathrm{gf}}_{\mathbf Q}(M_i)\) of \emph{rational}
generically finite Fubini--Study functions with \(c_1(M_i)\to\nu^*\theta\);
that is \eqref{eq:bj-positive-upstairs-regularization} verbatim, and no
perturbation of weights is needed.  That theorem is the Fubini--Study
refinement of \Cref{def:bj-psh}\,(3) and is proved in
\Cref{chap:bj-fs-pl-psh} with no envelope input.  It is \emph{not}
\Cref{cor:ch41-real-weights}, which runs the other way --- it extends the
determinant-norm averaging \(\operatorname{Av}_q\) to weights in an arbitrary
additive subgroup --- and invoking that corollary here would also contradict
the declaration above that none of the averaging machinery of
\Cref{chap:bj-finite-descent} is used.  Nothing in this step involves
\(X\).

\emph{The classes are promoted.}  \Cref{lem:bj-normalization-class-promotion}
replaces the presentations by presentations in classes \(\nu^*L_i\) with
\(c_1(L_i)\to\theta\), \emph{without changing any of the functions
\(\psi_i\)}.  Monotonicity of the net is therefore untouched, and so is the
common-zero locus of each stage.

\emph{Each stage descends.}  \Cref{prop:bj-frobenius-gf-descent} gives
\(\varphi_i\in\cH^{\mathrm{gf}}_{\mathbf Q}(L_i)\) with
\(\nu^*\varphi_i=\psi_i\).  Each \(\varphi_i\) lies in
\(\PSH_X(c_1(L_i))\): apply \Cref{def:bj-psh} with the constant net
\(\varphi_i\), which is decreasing, has the constant class \(c_1(L_i)\), and
is upper semicontinuous and generically finite by
\Cref{prop:ch34-generic-finiteness}\,\cref{itm:ch34-gf:cont} and
\cref{itm:ch34-gf:equiv}.

\emph{The descended net is decreasing with the right limit.}  Under the
identification \eqref{eq:bj-universal-analytic-homeomorphism} we have
\(\varphi_i=\psi_i\) as functions on the common space, so
\((\varphi_i)\) is decreasing and converges pointwise to the function
\(\varphi\) corresponding to \(\psi\).  This includes the points of proper
support and the points where the limit is \(-\infty\); by
\Cref{lem:ch43-identification-dictionary}\,\cref{itm:ch43-id:supp} nothing is
lost or gained.

\emph{The limit is psh.}  We have \(\varphi_i\in\PSH_X(c_1(L_i))\),
\(c_1(L_i)\to\theta\), \((\varphi_i)\) decreasing, and
\(\varphi(v_\triv)=\psi(w_\triv)\in\mathbf R\) by
\Cref{lem:ch43-identification-dictionary}\,\cref{itm:ch43-id:scaling} and
because \(\psi\) is psh, hence finite at the trivial valuation.  In
particular \(\varphi\not\equiv-\infty\), and \(X\) being integral this is
generic finiteness.  \Cref{thm:bj-basic-closure}\,(5) --- the
varying-class decreasing closure, that is, PSH2 --- gives
\(\varphi\in\PSH_X(\theta)\).  Finally \(\nu^*\varphi=\psi\) by
construction.  Note that only the \emph{decreasing} closure property was
used; no supremum was formed, and \Cref{warn:ch43-firewall} is respected.
\end{proof}

\begin{remark}[Where each hypothesis of
\Cref{thm:bj-positive-universal-reflection} is spent]
\label{rem:ch43-reflection-hypothesis-map}
\emph{Finiteness of \(\nu\)} enters twice: through
\Cref{lem:ch43-ample-pullback}, so that the buffer \(P_i\) can be taken
ample, and through the finite generation of \(B_U\) over \(A_U\) in
\Cref{lem:bj-uniform-frobenius}.  \emph{Universal homeomorphy} enters twice:
through \Cref{lem:bj-universal-analytic-homeomorphism}, and through the
elementwise criterion \Cref{lem:ch43-elementwise-frobenius}.
\emph{Characteristic \(p\)} enters exactly once, in
\Cref{lem:ch43-frobenius-subalgebra}, and everything downstream of it is
formal.  \emph{Projectivity} enters in the compactness argument of
\Cref{lem:bj-universal-analytic-homeomorphism}, in the existence of centers
used by \Cref{lem:ch43-zero-fs}, and in the availability of
\Cref{def:bj-psh}.  \emph{Integrality} of \(X\) and \(Y\) enters in the
injectivity of \(\cO_X\to\nu_*\cO_Y\) and in the reading of ``generically
finite''.  \emph{Birationality} enters nowhere
(\Cref{prop:ch43-frobenius-no-birational}).  \emph{Normality of \(Y\)}
enters nowhere in the proof; it is carried in
\eqref{eq:bj-universal-map} because the application is to a normalization
and because \Cref{chap:bj-envelope-consequences} will need
\Cref{thm:bj-normal-envelope} on \(Y\), which does require normality.
\end{remark}

\begin{corollary}
\label{cor:ch43-reflection-structure}
Assume \(\operatorname{char}k=p>0\).  The bijection
\eqref{eq:bj-positive-universal-reflection} is an isomorphism of ordered
sets; it commutes with the addition of real constants, with finite maxima,
with the scaling action of \(\mathbf R_{>0}\), and with decreasing limits
whose common pointwise limit is generically finite;
it carries \(\PSH_X(\theta)\cap\{\text{homogeneous}\}\) onto
\(\PSH_Y(\nu^*\theta)\cap\{\text{homogeneous}\}\); and it satisfies
\[
 \sup_{X^{\An}}\varphi=\varphi(v_\triv)
 =(\nu^*\varphi)(w_\triv)=\sup_{Y^{\An}}\nu^*\varphi .
\]
Moreover \(\PSH_X(\theta)\ne\emptyset\) if and only if
\(\PSH_Y(\nu^*\theta)\ne\emptyset\).
\end{corollary}

\begin{proof}
All the listed operations are pointwise, and
\eqref{eq:bj-positive-universal-reflection} is composition with the
homeomorphism \(\nu^{\An}\), which is equivariant for the scaling action by
\Cref{lem:ch43-identification-dictionary}\,\cref{itm:ch43-id:scaling} and
carries \(w_\triv\) to \(v_\triv\).  Homogeneity is the statement that
\(t\cdot\varphi=\varphi\) for the scaling action
(\Cref{def:homogeneous-function}), hence is preserved and reflected.  The
identity of suprema is \Cref{prop:ch32-maximum-principle}\,\cref{itm:ch32-mp:sup}
applied on each side: a psh function is decreasing for
\eqref{eq:semivaluation-order}, each of its Fubini--Study approximants being
so by
\Cref{prop:ch34-generic-finiteness}\,\cref{itm:ch34-gf:decreasing} and a
decreasing limit of decreasing functions being decreasing, so each supremum
is attained at the trivial valuation.  One may not cite
\Cref{prop:ch35-homogeneous-normalization} here: it is a statement about
\emph{homogeneous} psh functions, and applied to a general \(\theta\)-psh
\(\varphi\) it would assert \(\varphi(v_{\triv})=0\), which fails for
\(\varphi+\)\,constant.  The last assertion is immediate from
bijectivity.  Notice that \Cref{lem:ch34-psh-nonempty} only gives
\(\PSH_X(\theta)\ne\emptyset\Rightarrow\theta\) pseudoeffective; it does
not by itself identify nonemptiness with pseudoeffectivity in this
generality.
\end{proof}

\begin{warning}{Nef and bounded are not transported for free}
\label{warn:ch43-nef-not-transported}
\Cref{cor:ch43-reflection-structure} transports nonemptiness of the psh
class.  It does \emph{not} assert that
\(\theta\) is nef if and only if \(\nu^*\theta\) is, and the reader should
not read such a statement into it.  Nefness corresponds to the existence of
a \emph{bounded} psh function, and the reflection theorem says nothing about
boundedness beyond what is pointwise obvious.  In fact, since
\(\nu^{\An}\) is a homeomorphism, a function is bounded upstairs if and only
if its counterpart is bounded downstairs, so in this particular case the two
conditions do correspond; but the reason is the homeomorphism, not the
reflection theorem, and for a general finite morphism the corresponding
statement would require a separate argument.  We record the correct
statement: \(\theta\) is nef on \(X\) if and only if \(\PSH_X(\theta)\)
contains a bounded element, if and only if
\(\PSH_Y(\nu^*\theta)\) contains a bounded element, if and only if
\(\nu^*\theta\) is nef on \(Y\).  All three equivalences are pointwise
consequences of \Cref{lem:bj-universal-analytic-homeomorphism} and
\Cref{thm:bj-positive-universal-reflection}, and none of them uses the
envelope property.
\end{warning}

\begin{corollary}
\label{cor:ch43-ep-transport}
Assume \(\operatorname{char}k=p>0\) and let \(\theta\in N^1(X)_{\mathbf R}\).
\emph{If} \(\operatorname{EP}_Y(\nu^*\theta)\) holds --- that is, if every
nonempty bounded-above family in \(\PSH_Y(\nu^*\theta)\) has regularized
supremum in \(\PSH_Y(\nu^*\theta)\) --- then
\(\operatorname{EP}_X(\theta)\) holds.
\end{corollary}

\begin{proof}
Let \((\varphi_\alpha)\subseteq\PSH_X(\theta)\) be nonempty and bounded
above, and set \(u:=(\sup_\alpha\varphi_\alpha)^*\).  By
\Cref{cor:ch43-sup-star},
\[
 \nu^*u=\Bigl(\sup_\alpha\nu^*\varphi_\alpha\Bigr)^* .
\]
The family \((\nu^*\varphi_\alpha)\) lies in \(\PSH_Y(\nu^*\theta)\) by
\Cref{prop:ch43-free-half} and is bounded above by the same bound, since
\(\nu^{\An}\) is surjective.  The hypothesis
\(\operatorname{EP}_Y(\nu^*\theta)\) gives
\(\nu^*u\in\PSH_Y(\nu^*\theta)\).  Now
\Cref{thm:bj-positive-universal-reflection} supplies
\(u'\in\PSH_X(\theta)\) with \(\nu^*u'=\nu^*u\), and injectivity of
\(\nu^*\) --- which is injectivity of composition with a bijection --- gives
\(u=u'\in\PSH_X(\theta)\).
\end{proof}

\begin{remark}[This corollary is conditional, and deliberately so]
\label{rem:ch43-ep-conditional}
\Cref{cor:ch43-ep-transport} is the only place in this chapter where the
envelope property is mentioned, and it appears as a hypothesis on \(Y\), in
accordance with \Cref{warn:ch43-firewall}.  It is not combined here with the
normal case \Cref{thm:bj-normal-envelope}, although the combination is
immediate when \(Y=X^\nu\) is the normalization of a unibranch \(X\); that
assembly, together with its characteristic-zero half from
\Cref{chap:bj-subintegral-normalization}, is carried out in
\Cref{chap:bj-envelope-consequences}, where the resulting statement is
\Cref{thm:bj-final-envelope}.  Keeping the assembly in one place makes the
dependency diagram \eqref{eq:bj-final-dependency-diagram} auditable and
prevents the same theorem from being proved twice under two names.
\end{remark}

\begin{corollary}
\label{cor:ch43-reflection-tower}
Assume \(\operatorname{char}k=p>0\) and let
\[
 Z\xrightarrow{\ \nu_2\ }Y\xrightarrow{\ \nu_1\ }X
\]
be finite universal homeomorphisms of integral projective \(k\)-varieties.
Then \(\nu_1\circ\nu_2\) is one as well, and for every
\(\theta\in N^1(X)_{\mathbf R}\),
\[
 (\nu_1\nu_2)^*=\nu_2^*\circ\nu_1^*:
 \PSH_X(\theta)\xrightarrow{\ \sim\ }\PSH_Z\bigl((\nu_1\nu_2)^*\theta\bigr)
\]
is a bijection, and each factor is.  Moreover
\(e_{\min}(\nu_1\nu_2)\le e_{\min}(\nu_1)+e_{\min}(\nu_2)\).
\end{corollary}

\begin{proof}
That the composite is a finite universal homeomorphism is
\Cref{lem:ch43-uh-stability}\,\cref{itm:ch43-stab:comp}, together with the
fact that a composite of finite morphisms is finite.  Each factor gives a
bijection by \Cref{thm:bj-positive-universal-reflection}, applied with
\Cref{prop:ch43-frobenius-no-birational} so that birationality of the
factors need not be assumed, and the composite of bijections is a bijection;
the identity \((\nu_1\nu_2)^*=\nu_2^*\nu_1^*\) is functoriality of
composition with maps of sets.  For the exponents, write \(q_i=p^{e_i}\) with
\(e_i=e_{\min}(\nu_i)\), let \(U\subseteq X\) be affine, and let
\(c\in\cO_Z(( \nu_1\nu_2)^{-1}U)\).  Then \(c^{q_2}\in\cO_Y(\nu_1^{-1}U)\)
and hence \((c^{q_2})^{q_1}=c^{q_1q_2}\in\cO_X(U)\), so
\(e_1+e_2\in E(\nu_1\nu_2)\).
\end{proof}

\begin{remark}[On the bound for exponents in a tower]
\label{rem:ch43-tower-exponent-strict}
The bound is sharp.  Take \(X=Y=Z=\mathbf P^1\) and
\(\nu_1=\nu_2=F\) the Frobenius of \Cref{ex:ch43-frobenius-uh}, so that
\(e_{\min}(\nu_i)=1\).  The composite is \(F^2\), corresponding to the
inclusion \(k[\tau^{p^2}]\subset k[\tau]\), for which \(\tau^{p^e}\) lies in
\(k[\tau^{p^2}]\) if and only if \(p^2\mid p^e\), that is, \(e\ge2\).  Hence
\(e_{\min}(\nu_1\nu_2)=2=e_1+e_2\).

Equality is not the general rule.  Take \(\nu_2=\mathrm{id}\) and
\(\nu_1\) the normalization of the cuspidal cubic: then
\(e_{\min}(\nu_1\nu_2)=1=1+0\), trivially.  A less trivial phenomenon is
that the two inseparabilities may overlap, in which case the composite needs
fewer than \(e_1+e_2\) steps; only the inequality is proved above, and only
the inequality is used --- in fact it is used nowhere, and is recorded
because it makes the invariant \(e_{\min}\) behave like a length.
\end{remark}

\begin{corollary}\label{cor:bj-positive-normalization-reflection}
Let \(X\) be an integral projective unibranch variety over \(k\) with
\(\operatorname{char}k=p>0\), and let
\(\nu:X^\nu\to X\) be the normalization.
Then
\[
 \nu^*:\PSH_X(\theta)
       \xrightarrow{\ \sim\ }
       \PSH_{X^\nu}(\nu^*\theta)
\tag{43.18}\label{eq:bj-positive-normalization-reflection}
\]
for every real numerical class \(\theta\).
\end{corollary}

\begin{proof}
The normalization map satisfies the hypotheses of
\Cref{thm:bj-positive-universal-reflection} by
\Cref{lem:bj-unibranch-normalization-universal} below.
\end{proof}

\section{The unibranch normalization}
\label{sec:ch43-unibranch}

\subsection{Branches}
\label{sec:ch43-branches}

The hypothesis under which the envelope theorem is true is unibranchness,
and it is worth fixing exactly what the word means before the lemma that
converts it into a universal homeomorphism.

\begin{definition}[Unibranch]
\label{def:ch43-unibranch}
A local ring \(\cO\) is \emph{unibranch} if \(\cO_{\mathrm{red}}\) is a
domain whose integral closure in its fraction field is again local.  A
scheme is unibranch if all of its local rings are.  For an integral
noetherian scheme \(X\) with finite normalization \(\nu:X^\nu\to X\) this
says exactly that \(\nu^{-1}(x)\) is a single point for every \(x\in X\),
that is, that \(\nu\) is bijective; the number \(\#\nu^{-1}(x)\) is the
number of \emph{branches} of \(X\) at \(x\).
\end{definition}

The equality of the two descriptions is standard: the local ring of
\(X^\nu\) at a point above \(x\) is a localization of the integral closure
of \(\cO_{X,x}\), and the points above \(x\) correspond to the maximal
ideals of that integral closure, which is semilocal.  We shall use only the
second description --- \(\nu\) is bijective --- and it may be taken as the
definition for the purposes of this book, as it is in
\Cref{app:valuation-algebra}.

Two remarks, both used below.  \emph{Unibranch is strictly weaker than
normal}: a normal variety has \(\nu=\mathrm{id}\), hence is unibranch, and
the cuspidal cubic of \Cref{ex:ch43-cusp-basic} is unibranch and not normal.
\emph{Unibranch is strictly stronger than irreducible}: the nodal cubic is
integral and has two branches at its node
(\Cref{ex:ch43-node-not-unibranch}).

\begin{example}[The node is not unibranch]
\label{ex:ch43-node-not-unibranch}
Let \(N\subset\mathbf P^2\) be the nodal cubic \(\{ZY^2=X^2(X+Z)\}\), with
node at \(P=[0:0:1]\).  On the chart \(Z\ne0\), with \(x=X/Z\),
\(y=Y/Z\), the local ring at \(P\) is the localization of
\(k[x,y]/(y^2-x^2(x+1))\), whose normalization is \(k[t]\) with
\(x=t^2-1\), \(y=t(t^2-1)\); the two points \(t=\pm1\) both map to \(P\).
Hence \(\nu\) is not injective, \(N\) is not unibranch, and
\(\nu^{\An}\) is not injective either: the trivial valuations of the two
preimages are distinct points of \((\mathbf P^1)^{\An}\) with the same
image.  Everything in this chapter fails for \(N\), starting with
\Cref{lem:bj-universal-analytic-homeomorphism}; and indeed the envelope
property itself fails for an ample class on \(N\), which is why the
unibranch hypothesis is not removable.  See
\Cref{sec:ch43-failure} and \Cref{ex:bj-node-cusp}.
\end{example}

\begin{lemma}\label{lem:bj-unibranch-normalization-universal}
Let \(X\) be an integral projective unibranch variety over an algebraically
closed field.
Its normalization
\[
                         \nu:X^\nu\longrightarrow X
\tag{43.19}\label{eq:bj-unibranch-normalization}
\]
is a finite birational universal homeomorphism.
\end{lemma}

\begin{proof}
The scheme \(X\) is of finite type over a field and hence Nagata.
Its normalization is finite by
\cite[\href{https://stacks.math.columbia.edu/tag/035S}{Tag 035S}]
{stacks-project}.
It is birational by the definition of normalization of an integral scheme.
Unibranchness makes the normalization map bijective on the underlying
points.

We verify radiciality in the form needed later.
Because \(\nu\) is finite, its diagonal
\[
       \Delta:X^\nu\longrightarrow X^\nu\times_X X^\nu
\]
is a closed immersion.
If it were not surjective, its nonempty open complement would contain a
closed point.
The fibre product is of finite type over the algebraically closed field
\(k\), so that closed point would be \(k\)-rational.
It would determine two distinct closed points of \(X^\nu\) over one closed
point of \(X\), contradicting bijectivity.
Thus \(\Delta\) is surjective.
Tag 01S4 now shows that \(\nu\) is radicial.
Since \(\nu\) is finite, surjective, and radicial, Tag 04DF shows that it is a
universal homeomorphism.

The diagonal argument deserves to be written out, because it is the one
place in this chapter where a scheme-theoretic subtlety is genuinely used
and where an appeal to ``bijective, hence radicial'' would be a mistake
(\Cref{warn:ch43-bijective-not-enough}).

Let \(W:=X^\nu\times_XX^\nu\), a scheme of finite type over \(k\) because
\(X^\nu\) and \(X\) are.  The diagonal
\(\Delta:X^\nu\to W\) is a closed immersion because \(\nu\) is separated,
being finite; in particular \(\Delta(X^\nu)\) is a closed subset of \(W\)
and its complement \(V:=W\setminus\Delta(X^\nu)\) is open.  Suppose
\(V\ne\emptyset\).  A nonempty scheme of finite type over a field has a
closed point, and by the Nullstellensatz over the algebraically closed
\(k\) that closed point \(\omega\) has residue field \(k\).  A \(k\)-point
of \(W=X^\nu\times_XX^\nu\) is by definition a pair \((\eta_1,\eta_2)\) of
\(k\)-points of \(X^\nu\) with \(\nu(\eta_1)=\nu(\eta_2)\) as \(k\)-points
of \(X\); and \(\omega\notin\Delta(X^\nu)\) forces \(\eta_1\ne\eta_2\).
That contradicts the injectivity of \(\nu\) on closed points, which is
unibranchness.  Hence \(V=\emptyset\), \(\Delta\) is surjective, and
\Cref{prop:ch43-criteria}\,\cref{itm:ch43-crit:radicial} gives radiciality.
Surjectivity of \(\nu\) holds because \(\nu\) is finite and dominant, hence
closed with dense image.  Now
\Cref{prop:ch43-criteria}\,\cref{itm:ch43-crit:uh} applies.
\end{proof}

\begin{remark}[Where algebraic closedness enters]
\label{rem:ch43-where-algebraically-closed}
Bijectivity of a finite map on scheme points does not by itself force residue
field extensions to be trivial or purely inseparable over an arbitrary
ground field.
The closed-point argument above uses that every closed point of the finite
type fibre product is \(k\)-rational.
The theorem is stated over an algebraically closed field, exactly as required.

The counterexample of \Cref{warn:ch43-bijective-not-enough} shows that the
hypothesis is not decorative: over \(k=\mathbf F_p(s)\) the finite map
\(\Spec k(\sqrt s)\to\Spec k\), \(p\ne2\), is bijective on points and yet
not radicial; after base change to \(k(\sqrt s)\), the source has two points,
equivalently the diagonal is not surjective.
Over an algebraically closed field the phenomenon cannot occur, and the
proof above is exactly the observation that it cannot.
\end{remark}

\begin{corollary}
\label{cor:ch43-normalization-package}
Let \(X\) be an integral projective unibranch variety over the algebraically
closed field \(k\) and \(\nu:X^\nu\to X\) its normalization.  Then
\(\nu\) satisfies every hypothesis of \eqref{eq:bj-universal-map}: it is
finite, birational, a universal homeomorphism, and its source is a normal
integral projective variety.  Consequently
the sequence of results beginning with
\Cref{lem:bj-universal-analytic-homeomorphism} and ending with
\Cref{cor:ch43-ep-transport} applies to it verbatim, the last three under the additional hypothesis
\(\operatorname{char}k=p>0\).
\end{corollary}

\begin{proof}
Finiteness, birationality and universal homeomorphy are
\Cref{lem:bj-unibranch-normalization-universal}.  Normality of \(X^\nu\) is
the definition of the normalization; that it is integral follows since it
is the relative Spec of a sheaf of domains with fraction field \(K(X)\);
projectivity holds because \(\nu\) is finite, hence projective, and \(X\) is
projective.
\end{proof}

\begin{remark}[Unibranch is exactly the right hypothesis]
\label{rem:ch43-unibranch-sharp}
It is natural to ask whether unibranchness could be weakened.  It cannot,
and the reason is not a limitation of the method.  For an ample class on a
projective variety, the envelope property \emph{implies} unibranchness: this
is \Cref{thm:ep-implies-unibranch}, proved in
\Cref{chap:bj-envelope-history} from the results of Boucksom and Jonsson and
entirely independent of the present chapter.  The present chapter and
\Cref{chap:bj-subintegral-normalization} therefore prove the exact converse
of a known implication, and there is no room between the two.  The failure
on a non-unibranch variety can be seen directly on the nodal cubic
(\Cref{ex:ch43-node-not-unibranch}, \Cref{ex:bj-node-cusp}) and is
structural: with two branches, the analytic space upstairs is strictly
bigger than the one downstairs, and a psh function downstairs is
constrained by an equality between its values on the two branches which is
not preserved under taking suprema.
\end{remark}

\section{A nonnormal unibranch surface, end to end}
\label{sec:ch43-computations}

The three examples so far have been curves or have concerned a single step.
We now run the entire chapter --- analytic homeomorphism, class promotion,
uniform exponent, GF descent, decreasing closure --- on a nonnormal
unibranch surface, with a class that is genuinely irrational so that the
approximating classes must vary, and with values computed at divisorial and
at proper-support points.

\begin{example}[The cuspidal cylinder]
\label{ex:ch43-surface-endtoend}
Let \(k\) be algebraically closed with \(\operatorname{char}k=p>0\), let
\(C\subset\mathbf P^2\) be the cuspidal cubic of
\Cref{ex:ch43-cusp-basic}, and set
\[
 X:=C\times\mathbf P^1,
 \qquad
 Y:=\mathbf P^1\times\mathbf P^1,
 \qquad
 \nu:=\nu_C\times\mathrm{id}.
\]
Write \([u:w]\) and \([z_0:z_1]\) for the two sets of homogeneous
coordinates on \(Y\), and \(t:=u/w\), \(z:=z_0/z_1\).

\emph{(a) The hypotheses.}  \(X\) is an integral projective surface, singular
exactly along the curve \(\{P\}\times\mathbf P^1\), where \(P\) is the cusp.
The map \(\nu\) is finite (a product of a finite map with an identity),
birational, and bijective, so \(X\) is unibranch and
\Cref{lem:bj-unibranch-normalization-universal} makes \(\nu\) a finite
birational universal homeomorphism.  That \(Y\) is the normalization of
\(X\) follows because \(Y\) is normal and \(\nu\) is finite birational.  The
standing hypotheses \eqref{eq:bj-universal-map} are met.

\emph{(b) The charts and the exponent.}  Cover \(X\) by the four affine
opens \(U_\epsilon\times V_\delta\) with
\(U_1=C\cap\{Z\ne0\}\), \(U_2=C\cap\{Y\ne0\}\),
\(V_1=\{z_1\ne0\}\), \(V_2=\{z_0\ne0\}\).  On \(U_1\times V_1\),
\[
 A=k[t^2,t^3][z],
 \qquad
 B=k[t][z],
\]
and \(B=A[t]\) with \(t^p\in A\); the same computation applies verbatim on
\(U_1\times V_2\), with \(z\) replaced by \(z'=z_1/z_0\), because the
singular locus \(\{P\}\times\mathbf P^1\) meets both \(V_1\) and \(V_2\);
on the two charts over \(U_2\), \(\nu\) is an isomorphism.  By \Cref{lem:bj-uniform-frobenius} and
\Cref{lem:ch43-minimal-exponent},
\[
 q=p,
 \qquad
 e_{\min}(\nu)=1 .
\]
Note that the exponent is the same as for the curve \(C\): multiplying by a
smooth factor changes neither the generators nor the semigroup computation.

\emph{(c) The class.}  Let \(H_1:=\mathrm{pr}_1^*\cO_C(1)\) and
\(H_2:=\mathrm{pr}_2^*\cO_{\mathbf P^1}(1)\) on \(X\).  Identify
\(N^1(Y)_{\mathbf R}\) with \(\mathbf R^2\) by the bidegree, so that
\(c_1(\cO_Y(a,b))\leftrightarrow(a,b)\).  Then
\[
 \nu^*c_1(H_1)=(3,0),
 \qquad
 \nu^*c_1(H_2)=(0,1),
\]
the first because \(C\) has degree three and \(\nu_C\) is birational.  Put
\[
 \theta:=c_1(H_1)+\sqrt2\,c_1(H_2)\in N^1(X)_{\mathbf R},
 \qquad
 \nu^*\theta=(3,\sqrt2).
\]
The class is irrational, so no single rational line bundle carries it and
every regularization must have varying classes.

\emph{(d) A psh function upstairs.}  Set
\[
 \psi:=\log|uw^2|+\sqrt2\,\log|z_0| .
\]
Each summand is a GF function --- \(uw^2\in H^0(Y,\cO(3,0))\) and
\(z_0\in H^0(Y,\cO(0,1))\) --- so
\(\log|uw^2|\in\PSH_Y(3,0)\) and \(\log|z_0|\in\PSH_Y(0,1)\), and
\(\psi\in\PSH_Y(\nu^*\theta)\) by
\Cref{thm:bj-basic-closure}\,(3).  It is not a GF function, its class being
irrational.

\emph{(e) The regularization.}  Choose rationals \(r_n\uparrow\sqrt2\) and
write \(r_n=a_n/b_n\) in lowest terms.  Put
\[
 \psi_n:=\log|uw^2|+r_n\log|z_0|
 =\frac1{b_n}\log\bigl|(uw^2)^{\otimes b_n}\otimes z_0^{\otimes a_n}\bigr|
 \in\cH^{\mathrm{gf}}_{\mathbf Q}\bigl(Y,\cO(3,r_n)\bigr).
\]
Since \(\log|z_0|\le0\) everywhere --- a global section has nonnegative
value at every center --- the net \((\psi_n)\) is \emph{decreasing}, and it
converges pointwise to \(\psi\).  Its classes \((3,r_n)\) converge to
\(\nu^*\theta\).

\emph{(f) Promotion.}  Here the classes are already pullbacks:
\(\cO_Y(3,r_n)=\nu^*L_n\) with \(L_n:=H_1+r_nH_2\in\Pic(X)_{\mathbf Q}\), and
\(c_1(L_n)\to\theta\).  So \Cref{lem:bj-normalization-class-promotion} is
applied with the trivial buffer \(P_n=0\); the reader should not conclude
that the step is always vacuous.  It is vacuous here because
\(\nu^*:\Pic(X)_{\mathbf Q}\to\Pic(Y)_{\mathbf Q}\) happens to be
surjective, which is special to this \(X\).  When it is not surjective, or
when the given net is presented in bundles not of that form --- which
\Cref{def:bj-psh} in no way forbids --- the buffer is needed, and
\Cref{ex:ch43-promotion-numbers} shows what it looks like.

\emph{(g) Descent of the stage \(n\).}  On the chart \(U_1\times V_1\) take
the frame of \(b_nL_n=b_nH_1+a_nH_2\) given by \(Z^{\otimes b_n}\otimes
z_1^{\otimes a_n}\).  The section
\(\sigma_n:=(uw^2)^{\otimes b_n}\otimes z_0^{\otimes a_n}\) has coefficient
\[
 f=\frac{(uw^2)^{b_n}z_0^{a_n}}{w^{3b_n}z_1^{a_n}}
  =t^{b_n}z^{a_n}\in B=k[t][z],
\]
which for \(b_n=1\) does not lie in \(A\).  Raising to the power \(q=p\),
\[
 f^p=t^{pb_n}z^{pa_n}\in A=k[t^2,t^3][z],
\]
because \(pb_n\ge2\) lies in the semigroup \(\langle2,3\rangle\).  Writing
\(pb_n=2c+3d\) with \(c,d\ge0\) --- possible since \(pb_n\ge2\) --- and
noting \(c+d\le pb_n\), the glued section is
\[
 s_n^{[p]}
 =X^{c}Y^{d}Z^{\,pb_n-c-d}\otimes z_0^{\,pa_n}
 \in H^0\bigl(X,\ pb_nL_n\bigr),
\]
and the descended stage is
\[
 \varphi_n:=\frac1{pb_n}\log\bigl|s_n^{[p]}\bigr|
 \in\cH^{\mathrm{gf}}_{\mathbf Q}(X,L_n),
 \qquad
 \nu^*\varphi_n=\psi_n .
\]

\emph{(h) Values, computed.}  For \(\alpha,\beta\ge0\) let \(w_{\alpha\beta}\)
be the monomial semivaluation on \(Y\) with
\(w_{\alpha\beta}(t)=\alpha\), \(w_{\alpha\beta}(z)=\beta\), centered at the
point \(t=z=0\).  Then
\[
 \psi_n(w_{\alpha\beta})=-\alpha-r_n\beta,
 \qquad
 \psi(w_{\alpha\beta})=-\alpha-\sqrt2\,\beta,
\]
and downstairs, at the corresponding point \(v_{\alpha\beta}\) with
\(v(x)=2\alpha\), \(v(y)=3\alpha\), \(v(z)=\beta\),
\[
 \begin{aligned}
 \varphi_n(v_{\alpha\beta})
 &=\frac1{pb_n}\Bigl(-c\cdot2\alpha-d\cdot3\alpha-pa_n\beta\Bigr)\\
 &=\frac1{pb_n}\Bigl(-(2c+3d)\alpha-pa_n\beta\Bigr)\\
 &=\frac{-pb_n\alpha-pa_n\beta}{pb_n}
  =-\alpha-r_n\beta,
 \end{aligned}
\]
using \(2c+3d=pb_n\).  The two agree, as they must, and the answer does not
depend on the choice of \(c\) and \(d\) --- a reassuring check, since that
choice was arbitrary.  For
instance with \(p=5\), \(r_n=7/5\), \(\alpha=1\), \(\beta=2\):
\[
 \psi_n=-1-\tfrac{14}5=-\tfrac{19}5,
 \qquad
 \varphi_n=-\tfrac{19}5,
\]
while \(\psi=-1-2\sqrt2=-3.828\ldots\) and \(-19/5=-3.8\), consistent with
\(r_n\uparrow\sqrt2\) and \(\psi_n\downarrow\psi\).

\emph{(i) Proper support.}  Let \(v\) be the semivaluation of \(X\) whose
support is the singular curve \(\{P\}\times\mathbf P^1\); upstairs the unique
preimage is the semivaluation \(w\) with support \(\{t=0\}\).  Then
\(w(t)=+\infty\), so \(\psi_n(w)=-\infty\) and \(\psi(w)=-\infty\);
downstairs \(v(x)=v(y)=+\infty\), so the section \(s_n^{[p]}\) has
\(\log|s_n^{[p]}|(v)=-\infty\) as soon as \(c+d>0\), which holds because
\(2c+3d=pb_n\ge2\).  Hence \(\varphi_n(v)=-\infty\).  The equality
\(\nu^*\varphi_n=\psi_n\) therefore holds in \(\Rbar\) at this point too, and
the decreasing limit is \(-\infty\) there on both sides.  A proof restricted
to divisorial points would have lost exactly this information.

\emph{(j) The conclusion.}  \((\varphi_n)\) is a decreasing net in
\(\cH^{\mathrm{gf}}_{\mathbf Q}(X,L_n)\subseteq\PSH_X(c_1(L_n))\) with
\(c_1(L_n)\to\theta\) and \(\varphi_n(v_\triv)=0\) for every \(n\), so the
limit is not identically \(-\infty\).  By
\Cref{thm:bj-basic-closure}\,(5),
\[
 \varphi:=\lim_n\varphi_n\in\PSH_X(\theta),
 \qquad
 \nu^*\varphi=\psi .
\]
This is \Cref{thm:bj-positive-universal-reflection} for this \(\psi\), run
by hand.  Observe that at no point did we form a supremum, use the envelope
property, average over a fibre, or divide by a degree.
\end{example}

\begin{remark}[What the surface example adds to the curve examples]
\label{rem:ch43-surface-adds}
Three things.  First, the singular locus is positive dimensional, so the
loci at which values are \(-\infty\) are not isolated points and the
proper-support bookkeeping of item~(i) is not a degenerate case.  Second,
the class \(\theta\) is irrational, so the approximating classes genuinely
vary and \Cref{def:bj-psh} must be used in its varying-class form; a
formulation of the psh class with a fixed \(L\) would not even allow
\(\psi\) to be written down.  Third, the exponent \(q=p\) is computed on a
cover of four charts of which two are singular, both giving the same
\(e=1\), illustrating that the
maximum of \Cref{lem:bj-uniform-frobenius} is usually attained on a small
part of the cover, while \Cref{cor:ch43-frobenius-all-affines} is what allows
the cover to be chosen afterwards to trivialize \(b_nL_n\) --- and in
item~(g) the trivializing cover does depend on \(n\), through \(a_n\) and
\(b_n\).
\end{remark}

\section{Extended-real and circularity audit}
\label{sec:ch43-audit}

The positive-characteristic reflection theorem has three independent safety
checks.
First, \Cref{lem:bj-universal-analytic-homeomorphism} starts with the support
prime and therefore includes semivaluations with nonzero kernel.
Second, \eqref{eq:bj-positive-gf-identity} is obtained by powering all
coefficients and so remains literal at common zeros.
Third, PSH2 is applied only after the trivial valuation excludes the
identically \(-\infty\) alternative.

No envelope is formed anywhere in the proof.
The only psh inputs are pullback, GF regularization, and decreasing closure.
Therefore \Cref{thm:bj-positive-universal-reflection} is available before the
final envelope theorem and may safely be used in its proof.

\subsection{Where extended-real values occur, and what is done with them}
\label{sec:ch43-extended-real}

The chapter manipulates functions with values in \(\Rbar\), and the
conventions sheet forbids forming \((-\infty)-(-\infty)\).  The following is
a complete list of the places where an infinite value occurs, together with
the operation performed on it.

\begin{enumerate}
\item\label{itm:ch43-er:homeo}
      In \Cref{lem:bj-universal-analytic-homeomorphism} and
      \Cref{lem:ch43-identification-dictionary}, values do not occur at all:
      the statements are about points and about supports.  The
      \(-\infty\)-loci are transported because supports are, not because any
      value is manipulated.
\item\label{itm:ch43-er:usc}
      In \Cref{cor:bj-universal-usc} and \Cref{cor:ch43-sup-star}, suprema
      and infima are taken in \(\Rbar\).  These are order operations and are
      total; no subtraction occurs.
\item\label{itm:ch43-er:zero}
      In \Cref{lem:ch43-zero-fs} the conclusion is that a maximum equals
      \(0\), \emph{including} at proper-support points; the proof shows one
      term is \(0\) there, so the maximum is not \(-\infty\).  The lemma is
      used in \Cref{lem:bj-normalization-class-promotion} to \emph{add} a
      function known to be identically zero, which is legitimate for any
      value in \(\Rbar\).
\item\label{itm:ch43-er:split}
      In the proof of \Cref{lem:bj-normalization-class-promotion} the
      identity \(\max_{j,a}(A_j+B_a)=\max_jA_j+\max_aB_a\) is used with
      \(A_j\in[-\infty,\infty)\) and \(B_a\in\mathbf R\).  This is valid in
      \(\Rbar\): if some \(A_j=-\infty\) the corresponding terms are
      \(-\infty\) and are discarded by the maximum on both sides, and if all
      \(A_j=-\infty\) both sides are \(-\infty\).
\item\label{itm:ch43-er:power}
      In \Cref{prop:bj-frobenius-gf-descent} the identity
      \(w(f^q)=q\,w(f)\) is used with \(w(f)\in[0,+\infty]\).  It holds at
      \(+\infty\) because \(q\cdot(+\infty)=+\infty\); no cancellation is
      performed.
\item\label{itm:ch43-er:psh2}
      In \Cref{thm:bj-positive-universal-reflection} the decreasing limit is
      taken in \(\Rbar\), and the only thing that must be excluded is that
      it be identically \(-\infty\); this is excluded by evaluating at
      \(v_\triv\), where the value is finite.
\end{enumerate}

\begin{warning}{Three operations this chapter never performs}
\label{warn:ch43-forbidden}
\emph{No difference of two psh functions is formed}, and in particular no
expression of the shape \(\varphi-\nu^*\psi\) occurs; the reflection is
proved by exhibiting a presentation, not by estimating a difference.

\emph{No supremum of psh functions is formed}, regularized or not, except in
the hypothesis of \Cref{cor:ch43-ep-transport}, which is displayed as such.

\emph{No division by a fibre degree is performed.}  The only divisions are by
the integers \(m\), \(q\) and \(r_i\) and by the real number
\(r_im_i\); all of them are divisions in \(\mathbf R\) applied to a
real-valued or \(-\infty\)-valued maximum, never divisions in \(k\) or in a
finite \(k\)-algebra.  This is what makes the chapter insensitive to whether
\(p\) divides anything.
\end{warning}

\subsection{Checkpoints}
\label{sec:ch43-checkpoints}

A reader auditing the chapter may verify the following six checkpoints.

\emph{Checkpoint 1.}  \Cref{lem:bj-universal-analytic-homeomorphism} is
proved for all analytic points, not only for valuations of \(K(X)\); see
step three of its expanded proof and
\Cref{lem:ch43-identification-dictionary}\,\cref{itm:ch43-id:supp}.

\emph{Checkpoint 2.}  The exponent \(q\) is chosen before \(L\), \(m\) and
the trivializing cover; the licence for this is
\Cref{cor:ch43-frobenius-all-affines}, and the consequence of omitting it is
described in \Cref{rem:ch43-fixed-cover-not-enough}.

\emph{Checkpoint 3.}  The class promotion changes no value of any function;
see the explicit computation at the end of the proof of
\Cref{lem:bj-normalization-class-promotion} and
\Cref{rem:ch43-promotion-not-claimed}.

\emph{Checkpoint 4.}  The identity
\eqref{eq:bj-positive-gf-identity} is an identity in \(\Rbar\) at every
analytic point, verified at common zeros in the fifth paragraph of the
expanded proof.

\emph{Checkpoint 5.}  The only closure property of \(\PSH\) invoked is
\Cref{thm:bj-basic-closure}\,(5), for \emph{decreasing} nets with varying
classes.  See the last paragraph of the proof of
\Cref{thm:bj-positive-universal-reflection} and
\Cref{rem:ch43-firewall-not-vacuous}.

\emph{Checkpoint 6.}  No result of
\Cref{chap:bj-alterations,chap:bj-subintegral-normalization,%
chap:bj-envelope-consequences} is used in any proof.  The occurrences of
\Cref{thm:bj-normal-envelope} in this chapter are in
\Cref{rem:ch43-reflection-hypothesis-map} and
\Cref{rem:ch43-ep-conditional}, both of which are remarks about what happens
elsewhere.

\begin{proposition}\label{prop:bj-positive-normalization-ledger}
For the reflection theorem:
\begin{enumerate}
\item finiteness supplies finite affine algebras and a finite affine cover;
\item birationality makes the affine ring maps injective;
\item universal homeomorphy gives unique analytic lifts and the elementwise
      Frobenius-power criterion;
\item positive characteristic turns the elementwise criterion into actual
      coefficient descent by the Frobenius identity;
\item projectivity supplies compact Hausdorff analytifications and centered
      semivaluations;
\item normality is required only of the upstairs variety for the eventual
      normalization application and the available BJ regularization package.
\end{enumerate}
None of these hypotheses is replaced by the assertion that \(\nu\) is merely
bijective on closed points.
\end{proposition}

\begin{proof}
The first four items are exactly
\Cref{lem:bj-universal-analytic-homeomorphism,lem:bj-uniform-frobenius,prop:bj-frobenius-gf-descent}.
Projectivity was used in the compact-to-Hausdorff argument and in the centered
valuation reading of basepoint-free systems.
The last item records the setup of \eqref{eq:bj-universal-map}.

Item~(2) deserves a comment, since \Cref{prop:ch43-frobenius-no-birational}
shows that birationality can be dropped.  What is actually needed for the
affine ring maps to be injective is that \(\nu\) be dominant between
integral schemes, which follows from surjectivity; birationality is a
convenient sufficient condition and is present in the application, but is
not the operative hypothesis.  The precise map of hypotheses to uses is
\Cref{rem:ch43-reflection-hypothesis-map}.
\end{proof}

\section{Failure boundaries}
\label{sec:ch43-failure}

Four hypotheses have been used.  We indicate for each what goes wrong when
it is dropped, so that the theorem is not remembered in a stronger form than
it has.

\subsection{Dropping positive characteristic}
\label{sec:ch43-failure-char}

The topological half survives:
\Cref{lem:bj-universal-analytic-homeomorphism},
\Cref{cor:bj-universal-usc}, \Cref{cor:ch43-sup-star},
\Cref{lem:ch43-identification-dictionary},
\Cref{prop:ch43-free-half} and
\Cref{lem:bj-normalization-class-promotion} are proved in every
characteristic, and \Cref{lem:bj-unibranch-normalization-universal} likewise.
Everything from \Cref{sec:ch43-frobenius} onwards fails as stated: there is
no ring homomorphism \(x\mapsto x^q\), the set \(\{b:b^q\in A\}\) is not
closed under addition (\Cref{warn:ch43-subalgebra-charp}), and no uniform
exponent exists (\Cref{rem:ch43-charzero-contrast}).  The conclusion,
however, is still true, and is proved by a different mechanism in
\Cref{chap:bj-subintegral-normalization}: in characteristic zero a finite
universal homeomorphism is subintegral
(\Cref{rem:appA-subintegral-stronger}), and the Roberts--Singh additive-power
theorem \Cref{thm:appA-roberts-singh} replaces the Frobenius identity.  The
price is that the descended objects bracket the given function instead of
equalling it, so a limit and a monotonization are required.

\subsection{Dropping unibranchness}
\label{sec:ch43-failure-unibranch}

Then \(\nu\) is not injective, hence not radicial, hence not a universal
homeomorphism, and \Cref{lem:bj-universal-analytic-homeomorphism} fails at
its first step: distinct branches over a point give distinct analytic points
upstairs with the same image.  For the nodal cubic of
\Cref{ex:ch43-node-not-unibranch} the two preimages of the node give two
trivial semivaluations of \((\mathbf P^1)^{\An}\) mapping to the same point
of \(N^{\An}\), so \(\nu^{\An}\) is two-to-one there and \(\nu^*\) on
functions is no longer surjective onto anything: a function upstairs taking
different values at the two preimages is not a pullback.  This is not a
technical obstruction but the true one:
\Cref{thm:ep-implies-unibranch} shows that the envelope property for an
ample class \emph{forces} unibranchness, so no argument can remove the
hypothesis.  The book therefore does not place ``remove unibranchness'' in
any research queue.

\subsection{Dropping radiciality while keeping finiteness}
\label{sec:ch43-failure-radicial}

A finite surjective morphism that is not radicial is the subject of
\Cref{chap:bj-finite-descent}, and the mechanism there is completely
different: one averages over the fibre using a determinant norm, and the
descended function is the fibre average
\(\frac1d\sum_{w\mid z}e(w/z)\varphi(w)\), not \(\varphi\) itself.  In
particular pullback is then \emph{not} surjective onto the psh class
upstairs: its image consists of fibrewise constant functions, which is a
proper subset as soon as \(d>1\) and the map is not radicial
(\Cref{cor:ch41-finite-bijection}, \Cref{ex:ch42-not-a-pullback}).  What the
present chapter proves is a strictly stronger conclusion under a strictly
stronger hypothesis, and the two are not comparable as techniques.

\subsection{Dropping quasi-compactness or finite type}
\label{sec:ch43-failure-compact}

\Cref{ex:ch43-no-uniform-exponent} gives an affine universal homeomorphism
of \(\mathbf F_p\)-schemes admitting no uniform exponent.  It is not of
finite type, and by
\Cref{prop:ch43-criteria}\,\cref{itm:ch43-crit:finite} that is exactly the
loophole: a universal homeomorphism of finite type is finite, and then
\Cref{lem:bj-uniform-frobenius} applies.  Within the standing hypotheses of
this chapter --- \(X\) a projective variety --- the loophole is closed.  The
example is nonetheless worth keeping in mind when the theory is transported
to a nonprojective or nonnoetherian setting.

\subsection{What is not claimed about the exponent}
\label{sec:ch43-failure-exponent}

Finally, three statements that are \emph{not} proved and are not true in
general: that \(e_{\min}(\nu)\) is bounded in terms of \(\dim X\) or of the
degree of \(X\) in a fixed projective embedding
(\Cref{ex:ch43-no-uniform-exponent}); that \(e_{\min}(\nu)\) equals or
bounds the colength of the conductor
(\Cref{rem:ch43-emin-meaning}); and that the descended presentation can be
taken with the same denominator \(m\) as the given one
(\Cref{rem:ch43-not-presentations}).  Each of these would be a convenient
simplification and each is false.

\section{Chapter dictionary}
\label{sec:ch43-dictionary}

The correspondences established or used in this chapter are the following.
Rows one to three concern the morphism and its analytification, rows four
and five the algebra of the Frobenius exponent, and rows six to eight the
transport of Fubini--Study and psh functions.

\begin{center}
\begin{tabular}{p{.29\textwidth}p{.31\textwidth}p{.29\textwidth}}
\toprule
geometric datum & algebraic datum & hypothesis\\
\midrule
\(\nu\) a universal homeomorphism &
integral, surjective, universally injective; equivalently radicial with
surjective diagonal &
any base
(\Cref{def:bj-universal-homeomorphism}, \Cref{prop:ch43-criteria})\\
\(X\) unibranch &
\(\nu:X^\nu\to X\) bijective; diagonal of \(\nu\) surjective &
\(k=\bar k\), \(X\) integral of finite type
(\Cref{def:ch43-unibranch},
\Cref{lem:bj-unibranch-normalization-universal},
\Cref{warn:ch43-bijective-not-enough})\\
\(\nu^{\An}\) a homeomorphism &
unique extension \(|b|=|a|^{1/p^r}\) through a purely inseparable residue
extension &
\(\nu\) radicial surjective, \(X,Y\) projective
(\Cref{lem:ch43-inseparable-extension},
\Cref{lem:bj-universal-analytic-homeomorphism})\\
one Frobenius exponent per chart &
\(\{b:b^q\in A\}\) is an \(A\)-subalgebra; finitely many generators &
\(\operatorname{char}k=p>0\), \(\nu\) finite
(\Cref{lem:ch43-frobenius-subalgebra},
\Cref{lem:bj-uniform-frobenius})\\
one exponent for every affine open &
\(b^q\in A\) localizes at downstairs elements and glues &
\(X\) quasi-compact, \(\cO_X\hookrightarrow\nu_*\cO_Y\)
(\Cref{cor:ch43-frobenius-all-affines},
\Cref{rem:ch43-fixed-cover-not-enough})\\
class promotion &
\(M_i+P_i=\nu^*L_i\) with \(P_i\) ample and its complete linear system
contributing \(0\) &
\(\nu\) finite, \(Y\) projective
(\Cref{lem:ch43-zero-fs}, \Cref{lem:ch43-ample-pullback},
\Cref{lem:bj-normalization-class-promotion})\\
GF descent &
\(f_{\alpha j}^q\) glue to \(s_j^{[q]}\in H^0(X,qmL)\) with
\(\nu^*s_j^{[q]}=s_j^{\otimes q}\) &
\(\operatorname{char}k=p>0\), class already a pullback
(\Cref{prop:bj-frobenius-gf-descent}, \Cref{cor:ch43-gf-bijection})\\
psh reflection &
\(\nu^*:\PSH_X(\theta)\to\PSH_Y(\nu^*\theta)\) bijective, order and
finite maxima preserving &
\(\operatorname{char}k=p>0\); birationality not needed
(\Cref{thm:bj-positive-universal-reflection},
\Cref{prop:ch43-frobenius-no-birational})\\
\bottomrule
\end{tabular}
\end{center}

Four separations should be carried forward.

\emph{Homeomorphism versus isomorphism.}  \(\nu^{\An}\) is a homeomorphism
in every characteristic, and \(\nu\) is almost never an isomorphism.  The
first is a statement about points, the second about sheaves, and the whole
chapter lives in the gap between them
(\Cref{warn:ch43-homeo-not-iso}).

\emph{Exponent versus degree.}  The number \(q=p^e\) is an exponent
measuring an inseparability of structure sheaves; it is not the degree of
\(\nu\), which is one in the birational case, and it is never divided by.
The contrast with the number \(d\) of \Cref{chap:bj-finite-descent}, which
is a degree and \emph{is} divided by, is exact
(\Cref{rem:ch43-no-averaging}).

\emph{Elementwise versus uniform.}  The Stacks Project criterion is
elementwise; the descent argument needs a single exponent valid on every
affine open of \(X\).  The passage from one to the other consumes
characteristic \(p\) (for the subalgebra step), finiteness (for the
generators), quasi-compactness (for the maximum over charts) and the sheaf
axiom (for the localization).  Dropping any of the four breaks it
(\Cref{sec:ch43-uniformity}, \Cref{ex:ch43-no-uniform-exponent}).

\emph{Class versus bundle.}  A regularizing net delivers classes converging
to \(\nu^*\theta\); the Frobenius calculation needs bundles that are
literally pullbacks.  The gap is closed by adding an ample buffer whose
complete linear system contributes the zero function, so that the class
moves by \(o(1)\) and the function does not move at all
(\Cref{lem:bj-normalization-class-promotion},
\Cref{rem:ch43-numerical-insufficient}).

\section{Source and provenance}
\label{sec:ch43-source-note}

\Cref{lem:bj-universal-analytic-homeomorphism,lem:bj-normalization-class-promotion,lem:bj-uniform-frobenius,prop:bj-frobenius-gf-descent,thm:bj-positive-universal-reflection,lem:bj-unibranch-normalization-universal}
follow the positive-characteristic and unibranch portions of the
“Universal-homeomorphic normalization” section of the supplied envelope
manuscript.
Pullback and GF regularization are from \cite[Theorems~4.7 and~4.15]{BJ22}.
The Stacks Project inputs are exactly Tags 01S4, 04DF, 0CNF, and 035S, cited
at their points of use above.
The characteristic-zero statement suggested by the analytic homeomorphism is
not inferred formally:
its coefficient descent requires the separate additive-power argument of the
next chapter.

The manuscript is unpublished and has been made available to the author; its
results are quoted here as coming from a supplied manuscript, and none of
them is presented as a previously published theorem.  In its numbering, the
statements of this chapter correspond as follows.  Its Lemma on the
analytification of a universal homeomorphism is
\Cref{lem:bj-universal-analytic-homeomorphism}; its lemma promoting the
class of a regularizing net is
\Cref{lem:bj-normalization-class-promotion}; its lemma producing a single
Frobenius power valid for every open affine is the pair
\Cref{lem:bj-uniform-frobenius} and
\Cref{cor:ch43-frobenius-all-affines}; its proposition descending a GF
function in characteristic \(p\) is
\Cref{prop:bj-frobenius-gf-descent}; the characteristic-\(p\) half of its
theorem on descent through the normalization is
\Cref{thm:bj-positive-universal-reflection}; and its lemma identifying the
normalization of a unibranch variety is
\Cref{lem:bj-unibranch-normalization-universal}.  The characteristic-zero
half of the descent theorem, together with the additive-power sandwich on
which it rests, is the subject of
\Cref{chap:bj-subintegral-normalization} and is not duplicated here.

Two points of the exposition below are ours rather than the manuscript's,
and are flagged as such.  First, the manuscript states the uniform-exponent
lemma directly for every open affine \(U\subseteq X\), and gives the
localization argument in one sentence; because the distinction between ``a
fixed finite cover'' and ``every affine open'' is exactly what makes the
later frame calculation legal, we have separated the two statements and
given the gluing in full (\Cref{cor:ch43-frobenius-all-affines},
\Cref{rem:ch43-fixed-cover-not-enough}).  Second, the observation that
birationality is nowhere used in the characteristic-\(p\) argument
(\Cref{prop:ch43-frobenius-no-birational}) is an inspection of the proofs and
is not asserted in the manuscript, which carries birationality as a standing
hypothesis because its application is to a normalization; the observation is
harmless and is recorded because it clarifies which hypothesis does what.

The external attributions are as follows.  The psh formalism, the definition
of \(\cH^{\mathrm{gf}}_{\mathbf Q}\), the closure property under decreasing
nets with varying classes, and the regularization of psh functions by GF
functions are Boucksom--Jonsson \cite{BJ22}, with the corrections and
additions of \cite{BJ22b}; the book's versions are
\Cref{def:gf-fs-functions}, \Cref{def:bj-psh},
\Cref{thm:bj-basic-closure} and \Cref{thm:ch34-fs-regularization}, proved in
\Cref{chap:bj-fs-pl-psh}.  The four Stacks Project tags are audited in
\Cref{sec:stacks-universal-homeomorphisms} and
\Cref{sec:stacks-finite}, with the exact statements, the hypotheses to be
checked, and the first place of application; the dependency diagram is
\Cref{sec:stacks-dependency-chains}.  The necessity of unibranchness for an
ample class is \Cref{thm:ep-implies-unibranch}, from
\Cref{chap:bj-envelope-history}, and is independent of this chapter.  The
Roberts--Singh additive-power theorem, cited only for contrast, is
\cite{RobertsSingh}, imported as \Cref{thm:appA-roberts-singh}; the notion of
subintegrality goes back to Swan's work on seminormality, which we cite in
prose rather than by key.

Historically, the observation that a purely inseparable extension admits a
unique extension of an absolute value is as old as the valuation theory of
the 1930s, and appears in this form in the standard references; what is new
in the supplied manuscript is not that observation but the recognition that
one Frobenius power serves \emph{every} chart of \emph{every} cover at
once, so that the descent of a Fubini--Study function becomes a change of
Veronese multiple rather than an approximation procedure.  The reader who
wishes to see how much that recognition is worth should read
\Cref{sec:ch43-gf-descent} and then
\Cref{chap:bj-subintegral-normalization} in succession: the same conclusion
is reached, in the one case by an identity and in the other by a
quantitative sandwich, and the difference in length is the difference
between characteristic \(p\) and characteristic zero.
